\chapter{Translating environments}
\label{chap:trenv}

This chapter covers the environment layer of lean4lean's verification: the bridge between a
real \texttt{Lean.Kernel.Environment} (an \texttt{SMap Name ConstantInfo} together with a
quotient-initialization flag) and the abstract \texttt{VEnv} of the type theory (a partial map
to \texttt{VConstant}s, a set of defining equations, and, since the \texttt{iota} branch, a set
of registered pattern reduction rules \texttt{pats}). The central object is the inductive
relation $\mathtt{TrEnv}'\ safety\ C\ Q\ venv$ (\ref{ind:tr-env}), which has one constructor for
each way the kernel can extend an environment, and the top of the chapter is
\texttt{addDecl.WF} (\ref{thm:add-decl-wf}), the statement that a checked declaration addition
preserves that relation.

The files are
\srcloc{Lean4Lean/Verify/Environment/Basic.lean}{1} (619 lines: the translation of one constant,
the step relations \texttt{VEnv.AddConst}/\texttt{AddDef}/\texttt{AddQuot}, the
\texttt{insertConsts} map calculus, the inductive-block witness \texttt{AddInduct}, and
$\mathtt{TrEnv}'$),
\srcloc{Lean4Lean/Verify/Environment/Lemmas.lean}{1} (1178 lines: the \texttt{Aligned} companion
relation and its lookup calculus, $\delta$-soundness, the $\iota$ interface, structure-likeness,
a $\beta$-telescope calculus and projection reduction),
\srcloc{Lean4Lean/Verify/Environment/Quot.lean}{1} (615 lines: quotient initialization in full),
\srcloc{Lean4Lean/Verify/Environment/Checker.lean}{1} (230 lines: the per-declaration checks),
\srcloc{Lean4Lean/Verify/Environment/Extension.lean}{1} (487 lines: the reusable extension
machinery),
\srcloc{Lean4Lean/Verify/Environment.lean}{1} (208 lines: the top-level driver) and
\srcloc{Lean4Lean/Verify/Axioms.lean}{1} (510 lines: the project's trust boundary).

Thesis correspondences are indirect. Carneiro's thesis specifies the type theory, not the
kernel's environment data structure, so this group is a refinement layer with no counterpart of
its own; where the statements do touch the thesis they touch \texttt{axioms.tex}
\S\,``Definitions ($\delta$ reduction)'' (\texttt{TrConstant}, \texttt{TrEnv'.of\_value}),
\S\,``Inductive types'' and \S\,``The recursor'' (\texttt{TrIndType}, \texttt{TrRecursor},
\texttt{AddInduct}), \S\,``The computation rule ($\iota$ reduction)'' (\texttt{pats\_iota'} and
its inverse), \S\,``Quotient types'' (all of \texttt{Quot.lean}), and \texttt{typesys.tex}
\S\,\texttt{sec:undecidable}, whose $\mathsf{inv}_x$ construction is the shape of the projection
expansion that \texttt{proj\_defeq} reasons about.

Per-line attribution over the seven files is 1879 lines master, 1328 lines \texttt{iota} and 640
lines \texttt{trproj}. The \texttt{iota} contribution is concentrated in
\texttt{Basic.lean} lines 120--519 (the \texttt{insertConsts} calculus, \texttt{TrIndType},
\texttt{TrRecursor}, \texttt{AddInduct} and its derived lemmas), in the $\iota$ interface of
\texttt{Lemmas.lean}, and in the whole of \texttt{Quot.lean}. The \texttt{trproj} contribution is
\texttt{Lemmas.lean} lines 434--666, 771--912 and 949--1178 (structure-likeness,
\texttt{IotaRule} with the inverse lemma, the $\beta$ telescopes and \texttt{proj\_defeq}) plus
24 lines in \texttt{Basic.lean}, the parameters and fields it added to \texttt{TrIndType} and
\texttt{TrRecursor}.
\texttt{Checker.lean} and \texttt{Axioms.lean} are untouched by both branches. Sixteen of the
398 \texttt{iota} lines in \texttt{Basic.lean} (lines 34--49) are master code relocated by the
refactor commit \texttt{3046dff} rather than written on the branch.

Assessment. Within this group the \texttt{iota} branch is the stronger contribution: it replaces
a deliberately empty \texttt{AddInduct} (documented on master as ``essentially a \texttt{sorry}'')
with a real structure, derives its bookkeeping instead of assuming it, and is forced into, and
delivers, a full proof of quotient initialization that master had been getting for free from
vacuity; it also reaches the kernel model at one real point,
\texttt{inductiveReduceRecCore.WF}. The \texttt{trproj} contribution here is the more impressive
single theorem (\ref{thm:pats-iota-inv-shape} with \ref{thm:tr-env-proj-defeq}) and is correctly
architected, deriving the projection rule from the registered $\iota$ rule rather than
postulating it, but it is currently disconnected from the checker. Both rest on an undischarged
hypothesis: nothing in the repository constructs an \texttt{AddInduct} witness, and the
\texttt{inductDecl} case of \texttt{addDecl.WF} is still \texttt{sorry}
(\srcloc{Lean4Lean/Verify/Environment.lean}{208}), so the branches moved that \texttt{sorry}
rather than removing it.

Throughout, $C$ is a kernel constant map, $venv$ a model \texttt{VEnv}, $safety$ a
\texttt{DefinitionSafety}, and $env_1 \le env_2$ the model extension order.

\section{Translating one constant}

\begin{definition}[Safety level and value accessors of a \texttt{ConstantInfo}]
\label{family:constant-info-basics}
\lean{Lean.ConstantInfo.safety, Lean4Lean.ConstantInfo.hasValue_eq, Lean4Lean.ConstantInfo.value!_eq}
\leanok
\srcloc{Lean4Lean/Verify/Environment/Basic.lean}{11}
\texttt{ConstantInfo.safety} reads a declaration's \texttt{DefinitionSafety}: \texttt{unsafe} if
\texttt{isUnsafe}, else \texttt{partial} if \texttt{isPartial}, else \texttt{safe}. The two
accompanying one-line lemmas rewrite \texttt{hasValue} and \texttt{value!} into \texttt{value?}
(both by \texttt{cases ci}).
\end{definition}

\begin{definition}[Translation of a kernel constant, \texttt{TrConstant}]
\label{def:tr-constant}
\lean{Lean4Lean.TrConstant, Lean4Lean.TrConstVal, Lean4Lean.TrDefVal}
\leanok
\uses{ind:trexprs, family:constant-info-basics}
\srcloc{Lean4Lean/Verify/Environment/Basic.lean}{21}
$\mathtt{TrConstant}\ safety\ env\ ci\ ci'$ holds when $safety \le ci.\mathtt{safety}$, the number
of level parameters of $ci$ equals $ci'.\mathtt{uvars}$, and the kernel type of $ci$ translates by
\texttt{TrExprS} in the empty local context to $ci'.\mathtt{type}$. \texttt{TrConstVal} adds
$ci.\mathtt{name} = ci'.\mathtt{name}$; \texttt{TrDefVal} further requires
$ci.\mathtt{value!}$ (with \texttt{allowOpaque}) to translate to $ci'.\mathtt{value}$.
\thesisref{axioms.tex, \S ``Definitions ($\delta$ reduction)'': $\tau_{\bar\ell}(c)$ and
$v_{\bar\ell}(c)$}
\end{definition}

\begin{lemma}[Monotonicity of constant translation]
\label{family:tr-constant-mono}
\contrib{mixed}
\lean{Lean4Lean.TrConstant.sf_mono, Lean4Lean.TrConstant.mono, Lean4Lean.TrConstVal.mono, Lean4Lean.TrDefVal.mono}
\leanok
\srcloc{Lean4Lean/Verify/Environment/Basic.lean}{34}
\texttt{sf\_mono} lowers the safety index of a \texttt{TrConstant} along $safety \le safety'$; the
three \texttt{mono} lemmas transport a translation along an environment extension
$env \le env'$.
\end{lemma}
\begin{proof}
\uses{def:tr-constant, thm:trexprs-mono}
\leanok
Each unfolds the conjunction and applies \texttt{TrExprS.mono} (and, for \texttt{sf\_mono},
transitivity of the safety order).
\end{proof}

\begin{remark}
These four lemmas are marked as contributed by the per-line blame, but they are verbatim master
code relocated from \texttt{Lemmas.lean} into \texttt{Basic.lean} by the \texttt{iota} refactor
commit \texttt{3046dff}. The attribution is ``moved'', not ``written''.
\end{remark}

\begin{definition}[The model step of adding a constant or definition]
\label{def:venv-add-const}
\lean{Lean4Lean.VEnv.AddConst, Lean4Lean.VEnv.AddConst.le, Lean4Lean.VEnv.AddDef, Lean4Lean.VEnv.AddDef.le}
\leanok
\uses{def:tr-constant, def:venv-basic}
\srcloc{Lean4Lean/Verify/Environment/Basic.lean}{55}
$\mathtt{VEnv.AddConst}\ venv\ safety\ ci\ ci'\ venv'$ says: if $safety \le ci.\mathtt{safety}$
then $ci'$ translates $ci$, is well-formed, and $venv.\mathtt{addConst}\ ci.\mathtt{name}\ ci' =
\mathtt{some}\ venv'$; otherwise $venv' = venv$, matching the \texttt{ignore} step of
$\mathtt{TrEnv}'$. \texttt{VEnv.AddDef} is the analogue that adds the constant and then its
defining equation. Both \texttt{.le} lemmas give $venv \le venv'$. Stating this rather than just
$venv \le venv'$ is what lets a caller see which constant a step added.
\end{definition}

\begin{definition}[Refinement witness for quotient initialization, \texttt{AddQuot}]
\label{def:add-quot-witness}
\lean{Lean4Lean.AddQuot1, Lean4Lean.AddQuot, Lean4Lean.AddQuot1.to_addQuot, Lean4Lean.AddQuot1.le, Lean4Lean.AddQuot.to_addQuot, Lean4Lean.AddQuot.le}
\leanok
\uses{def:tr-constant, def:quot-consts, def:venv-addquot}
\srcloc{Lean4Lean/Verify/Environment/Basic.lean}{86}
$\mathtt{AddQuot1}\ name\ kind\ ci'\ P\ m\ env$ records one step of quotient initialization: for
some level parameters and type, the declaration
$\mathtt{quotInfo}\ \langle name, kind, \dots\rangle$ translates to the model constant $ci'$ at
safety \texttt{safe}, $name$ is fresh in $m$, $env.\mathtt{addConst}$ succeeds, and the
continuation $P$ holds of the extended map and environment. \texttt{AddQuot} chains four such
steps for \texttt{Quot}, \texttt{Quot.mk}, \texttt{Quot.lift}, \texttt{Quot.ind} and finishes with
\texttt{addDefEq quotDefEq}. \texttt{to\_addQuot} shows the model side is exactly
\texttt{VEnv.addQuot}, and \texttt{le} gives monotonicity.
\thesisref{axioms.tex, \S ``Quotient types''}
\end{definition}

\section{Inserting a block of constants}

\begin{definition}[Inserting a block into the constant map, \texttt{insertConsts}]
\label{def:insert-consts}
\contrib{iota}
\lean{Lean4Lean.insertConsts}
\leanok
\srcloc{Lean4Lean/Verify/Environment/Basic.lean}{127}
$\mathtt{insertConsts}\ C\ cis$ folds \texttt{SMap.insert} over $cis$, keying each
\texttt{ConstantInfo} by its own name. This is the map-level effect of adding a whole inductive
block, which the kernel performs one declaration at a time.
\thesisref{axioms.tex, \S ``Inductive types'' and \S ``The recursor''}
\end{definition}

\begin{lemma}[Lookup calculus for \texttt{insertConsts}]
\label{family:insert-consts-lemmas}
\contrib{iota}
\lean{Lean4Lean.insertConsts_cons, Lean4Lean.insertConsts_fresh_tail, Lean4Lean.insertConsts_wf, Lean4Lean.insertConsts_find?_mono, Lean4Lean.insertConsts_find?_mono_of_fresh, Lean4Lean.insertConsts_find?_none, Lean4Lean.insertConsts_find?, Lean4Lean.insertConsts_find?_self}
\leanok
\srcloc{Lean4Lean/Verify/Environment/Basic.lean}{130}
Under the hypotheses that $C$ is a well-formed map, that every name of the block is fresh in $C$
and that the block's names are pairwise distinct: the result is well-formed
(\texttt{insertConsts\_wf}); a lookup that succeeded before still succeeds with the same value
(\texttt{insertConsts\_find?\_mono}, \texttt{\_mono\_of\_fresh}); a lookup that succeeds after
either succeeded before or is one of the block's constants under its own name
(\texttt{insertConsts\_find?}); and each block constant resolves to itself
(\texttt{insertConsts\_find?\_self}).
\end{lemma}
\begin{proof}
\uses{def:insert-consts}
\leanok
All by induction on the block, with \texttt{insertConsts\_fresh\_tail} maintaining freshness of
the tail across the head insertion, and reducing at each step to the \texttt{SMap} lemmas
\texttt{SMap.find?\_insert}, \texttt{SMap.find?\_insert\_none} and
\texttt{SMap.insertList\_find?\_mono}.
\end{proof}

\reviewnote{\texttt{insertConsts\_find?\_none} (\srcloc{Lean4Lean/Verify/Environment/Basic.lean}{165})
is dead code: apart from its own statement, its only occurrence is its own recursive call.}

\section{The inductive-block witness}

\begin{definition}[Translation of one type former with its constructors, \texttt{TrIndType}]
\label{struct:tr-ind-type}
\contrib{mixed}
\lean{Lean4Lean.TrIndType}
\leanok
\uses{def:tr-constant, struct:vinductivetype, def:pi-telescope}
\srcloc{Lean4Lean/Verify/Environment/Basic.lean}{216}
$\mathtt{TrIndType}\ safety\ env_1\ env_T\ nparams\ block\ ival\ cvals\ t$ says: the kernel
\texttt{inductInfo} $ival$ translates to the model \texttt{VInductiveType} $t$ in the environment
$env_1$ before the block; $ival.\mathtt{ctors}$ is exactly the list of names of $cvals$; each
\texttt{ctorInfo} of $cvals$ translates in $env_T$, the environment already holding the block's
type formers, and satisfies
$cval.\mathtt{numParams} + cval.\mathtt{numFields} = c.\mathtt{type}.\mathtt{piArity}$;
$ival.\mathtt{all} = block$; $ival.\mathtt{numParams} = nparams$; and
$ival.\mathtt{numIndices} = t.\mathtt{type}.\mathtt{piArity} - nparams$.
\thesisref{axioms.tex, \S ``Inductive types'' and \S ``The recursor''}
\end{definition}

\begin{remark}
The core of \texttt{TrIndType} (the fields \texttt{tr}, \texttt{ctor\_names}, \texttt{ctors}) is
\texttt{iota}; the parameters $nparams$ and $block$ and the fields \texttt{all},
\texttt{numParams}, \texttt{numIndices} are \texttt{trproj} additions. The arity equation
$cval.\mathtt{numParams} + cval.\mathtt{numFields} = c.\mathtt{type}.\mathtt{piArity}$ is the
second conjunct of the \texttt{ctors} field and so is \texttt{iota} as well, not a \texttt{trproj}
addition; it is the bridge that makes \ref{thm:add-induct-ctor-find} and
\ref{thm:tr-env-ctor-arity} usable. Note
that the witness mentions three different environments: type formers translate in $env_1$,
constructors in $env_T$, and (in \ref{struct:tr-recursor}) reducts in $env_R$. That is correct,
but demanding on a producer.
\end{remark}

\begin{definition}[Translation of one recursor with its rules, \texttt{TrRecursor}]
\label{struct:tr-recursor}
\contrib{iota}
\lean{Lean4Lean.TrRecursor}
\leanok
\uses{def:tr-constant, struct:vrecursor, ind:trexprs}
\srcloc{Lean4Lean/Verify/Environment/Basic.lean}{237}
$\mathtt{TrRecursor}\ safety\ env_C\ env_R\ m_2\ rval\ r$ says the kernel \texttt{recInfo} $rval$
translates to the model \texttt{VRecursor} $r$ in $env_C$, the environment holding the type
formers and constructors; the telescope split (\texttt{numParams}, \texttt{numMotives},
\texttt{numMinors}, \texttt{numIndices}), the \texttt{all} list and the K-like flag \texttt{k} are
copied from $rval$; and the rules are matched one-to-one, each firing on the same constructor and
field count, with the constructor's parameter count read off its \texttt{ctorInfo} in the final
map $m_2$ and a reduct translating in $env_R$. The \texttt{trproj} field \texttt{name\_major}
adds that the kernel names a recursor $\mathtt{mkRecName}\ n$ after the type former $n$ its major
premise eliminates.
\thesisref{axioms.tex, \S ``The computation rule ($\iota$ reduction)''}
\reviewnote{The fields \texttt{all} (\srcloc{Lean4Lean/Verify/Environment/Basic.lean}{240}) and
\texttt{k} (\srcloc{Lean4Lean/Verify/Environment/Basic.lean}{245}) are never used anywhere in the
repository; \texttt{k} in particular records a flag for which the model has no rule. Looking up
the constructor in $m_2$ rather than in the block, by contrast, is a deliberate and documented
choice for nested blocks, whose auxiliary recursors fire on older constructors.}
\end{definition}

\begin{definition}[The constants a block adds, in model stage order]
\label{def:add-induct-consts}
\contrib{iota}
\lean{Lean4Lean.AddInduct.consts}
\leanok
\uses{family:ind-add-stages}
\srcloc{Lean4Lean/Verify/Environment/Basic.lean}{259}
$\mathtt{AddInduct.consts}\ ivals\ rvals$ is the list of \texttt{ConstantInfo}s a block
contributes: all type formers, then all constructors in block order, then all recursors. This
mirrors \texttt{VInductDecl.consts}, the order in which the model's \texttt{addTypesCtorsRecs}
folds \texttt{addConst}. The kernel uses this order for a non-nested block and inserts type by
type for a nested one, which is why \ref{struct:add-induct} records the actual insertion order
separately.
\end{definition}

\begin{definition}[Refinement witness for an inductive block, \texttt{AddInduct}]
\label{struct:add-induct}
\contrib{iota}
\lean{Lean4Lean.AddInduct}
\leanok
\uses{struct:tr-ind-type, struct:tr-recursor, def:add-induct-consts, def:insert-consts, family:ind-add-stages}
\srcloc{Lean4Lean/Verify/Environment/Basic.lean}{272}
$\mathtt{AddInduct}\ safety\ m_1\ env_1\ decl\ m_2\ env_2$ bundles the kernel-side data
($ivals$, $rvals$), the three intermediate model environments $env_T$, $env_C$, $env_R$, the four
successful stages $\mathtt{decl.addTypes}\ env_1 = \mathtt{some}\ env_T$,
$\mathtt{addCtors}\ env_T = \mathtt{some}\ env_C$,
$\mathtt{addRecs}\ env_C = \mathtt{some}\ env_R$,
$\mathtt{addRules}\ env_R = \mathtt{some}\ env_2$, a \texttt{TrIndType} for every type former and
a \texttt{TrRecursor} for every recursor, freshness in $m_1$ of every inserted name, and the
equation $m_2 = \mathtt{insertConsts}\ m_1\ order$ for an insertion order $order$ that is a
permutation of $\mathtt{consts}\ ivals\ rvals$.
\thesisref{axioms.tex, \S ``Inductive types'' and \S ``The recursor''}
\end{definition}

\begin{remark}
On master this was a deliberately empty \texttt{Prop} inductive whose docstring called it
``essentially a \texttt{sorry}'': it had no constructors, so the \texttt{induct} step of
$\mathtt{TrEnv}'$ could never fire and environments containing inductives were outside the
verified relation. Replacing it is the single largest structural change of the \texttt{iota}
branch. Two consequences deserve review. First, \texttt{AddInduct} is declared without a
\texttt{: Prop} ascription, so it is \texttt{Type}-valued and the \texttt{induct} constructor of
$\mathtt{TrEnv}'$ now carries data. This is needed because the derived lemmas project its fields
(\texttt{ivals}, \texttt{envR}, \texttt{order}), which a \texttt{Prop}-valued structure would not
allow. The often-stated consequence ``$\mathtt{TrEnv}'$ loses large elimination'' would be wrong:
$\mathtt{TrEnv}'$ has nine constructors, so it never had large elimination. What is true is that
an \texttt{AddInduct} witness cannot be recovered from a $\mathtt{TrEnv}'$ proof except into
\texttt{Prop}, which is all the present consumers need; the change to the shape of the project's
central invariant relative to master's \texttt{inductive AddInduct \ldots : Prop} is nowhere
documented. Second,
nothing in the repository ever constructs an \texttt{AddInduct}, so every result below is
conditional on a witness that \ref{thm:add-decl-wf} still \texttt{sorry}s.
\end{remark}

\begin{lemma}[Recomposing the four stages of \texttt{VEnv.addInduct}]
\label{family:add-induct-stages}
\contrib{iota}
\lean{Lean4Lean.AddInduct.env_eq, Lean4Lean.AddInduct.addTypesCtors, Lean4Lean.AddInduct.addTypesCtorsRecs, Lean4Lean.AddInduct.le, Lean4Lean.AddInduct.leR}
\leanok
\srcloc{Lean4Lean/Verify/Environment/Basic.lean}{296}
From the four stage equations of an \texttt{AddInduct} witness: \texttt{env\_eq} recomposes
$env_1.\mathtt{addInduct}\ decl = \mathtt{some}\ env_2$; \texttt{addTypesCtors} and
\texttt{addTypesCtorsRecs} recompose the intermediate prefixes; \texttt{le} gives
$env_1 \le env_2$ and \texttt{leR} gives $env_R \le env_2$.
\end{lemma}
\begin{proof}
\uses{struct:add-induct, def:ind-add-induct, family:indlem-stage-le}
\leanok
Each unfolds \texttt{VEnv.addInduct} (respectively its prefixes) and rewrites with the stored
stage equations; \texttt{leR} is \texttt{VEnv.addRules\_le} applied to the last stage.
\end{proof}

\begin{lemma}[Membership in a block's constants, by kind]
\label{family:add-induct-membership}
\contrib{iota}
\lean{Lean4Lean.AddInduct.mem_consts, Lean4Lean.AddInduct.novalue}
\leanok
\srcloc{Lean4Lean/Verify/Environment/Basic.lean}{312}
\texttt{mem\_consts} characterises $ci \in \mathtt{consts}\ ivals\ rvals$ three ways: $ci$ is a
type former, a constructor of some type former, or a recursor. \texttt{novalue} says every such
constant has $\mathtt{value?} = \mathtt{none}$ and $\mathtt{deltaValue?} = \mathtt{none}$, so a
block registers nothing $\delta$-reducible.
\end{lemma}
\begin{proof}
\uses{def:add-induct-consts}
\leanok
\texttt{mem\_consts} by \texttt{simp} on the append and \texttt{flatMap} and a case split in each
direction; \texttt{novalue} by \texttt{rcases} on \texttt{mem\_consts} followed by
\texttt{rfl} in each of the three cases.
\end{proof}

\begin{lemma}[The block's names are the model's, and are distinct]
\label{family:add-induct-names}
\contrib{iota}
\lean{Lean4Lean.AddInduct.types_names, Lean4Lean.AddInduct.ctors_names, Lean4Lean.AddInduct.recs_names, Lean4Lean.AddInduct.consts_names, Lean4Lean.AddInduct.names_nodup, Lean4Lean.AddInduct.order_fresh, Lean4Lean.AddInduct.order_nodup}
\leanok
\srcloc{Lean4Lean/Verify/Environment/Basic.lean}{334}
The kernel-side names of the type formers, constructors and recursors coincide with the model
declaration's (\texttt{types\_names}, \texttt{ctors\_names}, \texttt{recs\_names}, assembled by
\texttt{consts\_names}), and they are pairwise distinct (\texttt{names\_nodup}).
\texttt{order\_fresh} and \texttt{order\_nodup} transport freshness and distinctness to the
kernel's actual insertion order through the permutation.
\end{lemma}
\begin{proof}
\uses{struct:add-induct, family:indlem-find-fresh-nodup}
\leanok
The name equalities are the pointwise \texttt{map\_eq} lemma for \texttt{List.Forall}$_2$ on the \texttt{tr} fields.
\texttt{names\_nodup} splits the three-way append and uses that each model stage's
\texttt{addConst} fold succeeded (\texttt{VEnv.addTypes\_nodup}, \texttt{addCtors\_nodup},
\texttt{addRecs\_nodup}) together with cross-stage freshness (\texttt{addCtors\_fresh},
\texttt{addRecs\_fresh}); distinctness is derived, not assumed as a field.
\end{proof}

\begin{lemma}[An inductive block preserves constant-map well-formedness]
\label{thm:add-induct-wf}
\contrib{iota}
\lean{Lean4Lean.AddInduct.wf}
\leanok
\srcloc{Lean4Lean/Verify/Environment/Basic.lean}{395}
If $m_1$ is a well-formed \texttt{SMap} then so is $m_2$.
\end{lemma}
\begin{proof}
\uses{family:insert-consts-lemmas, family:add-induct-names, struct:add-induct}
\leanok
Rewrite by the \texttt{map\_eq} field and apply \texttt{insertConsts\_wf} with
\texttt{order\_fresh} and \texttt{order\_nodup}.
\end{proof}

\begin{lemma}[Lookup transport across an inductive block]
\label{family:add-induct-find}
\contrib{iota}
\lean{Lean4Lean.AddInduct.find?_mono, Lean4Lean.AddInduct.find?, Lean4Lean.AddInduct.find?_self, Lean4Lean.AddInduct.value_find}
\leanok
\srcloc{Lean4Lean/Verify/Environment/Basic.lean}{400}
A constant resolvable in $m_1$ is resolvable in $m_2$ to the same value
(\texttt{find?\_mono}); a constant resolvable in $m_2$ was either resolvable in $m_1$ or is one of
the block's constants under its own name (\texttt{find?}); each block constant resolves to itself
(\texttt{find?\_self}); and a constant that has a $\delta$-value and is resolvable in $m_2$ was
already resolvable in $m_1$ (\texttt{value\_find}).
\end{lemma}
\begin{proof}
\uses{family:insert-consts-lemmas, family:add-induct-membership, family:add-induct-names, struct:add-induct}
\leanok
All by the \texttt{map\_eq} field plus the corresponding \texttt{insertConsts} lemma, with the
permutation moving between $order$ and $\mathtt{consts}$; \texttt{value\_find} additionally uses
\texttt{novalue} to rule out the block case.
\end{proof}

\begin{theorem}[Kernel recursor lookup, resolved against the block]
\label{thm:add-induct-rec-find}
\contrib{iota}
\lean{Lean4Lean.AddInduct.rec_find}
\leanok
\uses{struct:add-induct, struct:tr-recursor}
\srcloc{Lean4Lean/Verify/Environment/Basic.lean}{434}
Assume $m_1$ is a well-formed \texttt{SMap}. If $m_2$ resolves $recName$ to a kernel
\texttt{recInfo} $rval$, then either $m_1$ already did, or
there is $r \in decl.\mathtt{recs}$ with $r.\mathtt{name} = recName$, the same major index and the
same telescope split, such that every kernel rule of $rval$ has a model rule with the same
constructor, the same field count, the constructor's parameter count as recorded in $m_2$, a
closed reduct, and a reduct translating the kernel rule's \texttt{rhs} in $env_2$.
\thesisref{axioms.tex, \S ``The computation rule ($\iota$ reduction)''}
\end{theorem}
\begin{proof}
\uses{family:add-induct-find, family:indlem-closed, family:add-induct-stages}
\leanok
\texttt{AddInduct.find?} splits the lookup; in the block case
\texttt{recs.forall\_exists\_l} produces the \texttt{TrRecursor}, the telescope-split fields give
the equality of \texttt{getMajorIdx}, and \texttt{VEnv.addRules\_closed} together with the
$env_R \le env_2$ monotonicity of \texttt{leR} supply closedness and the translation in $env_2$.
\end{proof}

\reviewnote{The conclusion is a long existential conjunction rather than a named structure,
unlike the dual \ref{thm:add-induct-rec-reg} which feeds \ref{struct:iota-rule}. The asymmetry in
style makes the two halves of the $\iota$ interface harder to compare than they need to be.}

\begin{theorem}[Model recursor is registered in the kernel map]
\label{thm:add-induct-rec-reg}
\contrib{iota}
\lean{Lean4Lean.AddInduct.rec_reg}
\leanok
\uses{struct:add-induct, struct:tr-recursor}
\srcloc{Lean4Lean/Verify/Environment/Basic.lean}{464}
Assume $m_1$ well-formed and, additionally, $decl.\mathtt{WF}\ env_1$. For each
$r \in decl.\mathtt{recs}$ there is a
kernel \texttt{RecursorVal} $rval$ with
$m_2.\mathtt{find?}\ r.\mathtt{name} = \mathtt{some}\ (\mathtt{recInfo}\ rval)$, the same major
index and telescope split, and such that every model rule of $r$ is found, keyed by constructor,
among $rval$'s rules by \texttt{List.find?}, with the same field count, the constructor's
parameter count in $m_2$, and a reduct translating the kernel rule's.
\thesisref{axioms.tex, \S ``The computation rule ($\iota$ reduction)''}
\end{theorem}
\begin{proof}
\uses{family:add-induct-find, structure:ind-wf, family:add-induct-stages}
\leanok
\texttt{recs.forall\_exists\_r} gives the \texttt{TrRecursor}; \texttt{find?\_self} places the
recursor in $m_2$ under its name; and \texttt{List.find?\_eq\_of\_nodup\_map} turns rule
membership into a \texttt{List.find?} result, using distinctness of the kernel rules' constructors
transported from \texttt{VInductDecl.WF.rules\_nodup}. The hypothesis $decl.\mathtt{WF}$ is used
only for that distinctness.
\end{proof}

\begin{theorem}[Kernel constructor lookup yields its $\iota$ rule]
\label{thm:add-induct-ctor-find}
\contrib{iota}
\lean{Lean4Lean.AddInduct.ctor_find}
\leanok
\uses{struct:add-induct, struct:tr-ind-type}
\srcloc{Lean4Lean/Verify/Environment/Basic.lean}{494}
Assume $m_1$ well-formed and $decl.\mathtt{WF}\ env_1$. If $m_2$ resolves $ctorName$ to a kernel \texttt{ctorInfo}
$cval$, then either $m_1$ already did, or there are $t \in decl.\mathtt{types}$ and
$c \in t.\mathtt{ctors}$ named $ctorName$, and a recursor $r \in decl.\mathtt{recs}$ with a rule
$ru$ firing on $ctorName$ such that $ru.\mathtt{ctorParams} = cval.\mathtt{numParams}$,
$ru.\mathtt{nfields} = cval.\mathtt{numFields}$ and $r.\mathtt{numParams} = cval.\mathtt{numParams}$.
\thesisref{axioms.tex, \S ``The computation rule ($\iota$ reduction)''}
\end{theorem}
\begin{proof}
\uses{family:add-induct-find, family:ind-wf-corollaries, thm:indlem-rules-ctor-shape}
\leanok
\texttt{AddInduct.find?} splits; then the \texttt{types} and \texttt{ctors} components of the
witness give the \texttt{TrIndType} data, \texttt{VInductDecl.WF.ctors\_have\_rules} produces the
rule, \texttt{rules\_own\_params} identifies \texttt{ctorParams} with $r.\mathtt{numParams}$, and
\texttt{rules\_ctor\_shape} together with the arity field of \texttt{TrIndType.ctors} forces
$ru.\mathtt{nfields} = cval.\mathtt{numFields}$ by \texttt{omega}.
\end{proof}

\begin{remark}
The chapter's source history shows this lemma, and the arity conjunct of
\texttt{TrIndType.ctors} it rests on, entering on \texttt{iota} in the commit
\texttt{d69ac5d} ``derive \texttt{AddInduct}'s $\iota$ bookkeeping instead of assuming it'', not on
\texttt{trproj}; per-line blame agrees. It is the only place the arity conjunct of
\texttt{TrIndType} earns its keep, and its only consumer is
\ref{thm:tr-env-structure-rec}, that is, the projection work: the $\iota$ pipeline itself never
calls it.
\end{remark}

\section{The translation relation}

\begin{definition}[Mutual definition blocks: map insertion and block translation]
\label{def:insert-defs}
\lean{Lean4Lean.insertDefs, Lean4Lean.TrDefBlock}
\leanok
\uses{def:tr-constant}
\srcloc{Lean4Lean/Verify/Environment/Basic.lean}{522}
$\mathtt{insertDefs}\ C\ cis$ inserts a list of \texttt{DefinitionVal}s as \texttt{defnInfo}s.
$\mathtt{TrDefBlock}\ safety\ env\ env'\ cis\ cis'$ pairs each kernel definition with a
\texttt{VDefVal}, translating headers against the environment before the block and values against
$env'$, the environment already holding every constant of the block; this mirrors the kernel
adding the headers as axioms first.
\thesisref{axioms.tex, \S ``Definitions ($\delta$ reduction)''}
\end{definition}

\begin{definition}[The environment translation relation, \texttt{TrEnv}]
\label{ind:tr-env}
\contrib{mixed}
\lean{Lean4Lean.TrEnv', Lean4Lean.TrEnv}
\leanok
\uses{struct:add-induct, def:add-quot-witness, def:insert-defs, def:tr-constant}
\srcloc{Lean4Lean/Verify/Environment/Basic.lean}{535}
$\mathtt{TrEnv}'\ safety\ C\ Q\ venv$ is an inductive relation between a kernel constant map $C$,
a quotient-initialization flag $Q$ and a model environment $venv$, generated by nine steps:
\texttt{empty}; \texttt{ignore} (a declaration not visible at $safety$, which extends the map
only); \texttt{axiom}; \texttt{defn} (constant plus defining equation); \texttt{mutualDef};
\texttt{thm} (whose model constant must be a \texttt{Prop}); \texttt{opaque} (no defining
equation); \texttt{quot} (an \texttt{AddQuot} witness, from $Q = \mathtt{false}$ to
$Q = \mathtt{true}$); and \texttt{induct} (an \texttt{AddInduct} witness with $decl.\mathtt{WF}$).
$\mathtt{TrEnv}\ safety\ env\ venv$ instantiates it at $env.\mathtt{constants}$ and
$env.\mathtt{quotInit}$.
\thesisref{axioms.tex, the admissible-declaration judgments of \S\S ``Definitions'', ``Inductive
types'' and ``Quotient types''}
\end{definition}

\begin{remark}
This is a master relation; the \texttt{iota} branch changed exactly one line, the
\texttt{AddInduct} argument of the \texttt{induct} constructor
(\srcloc{Lean4Lean/Verify/Environment/Basic.lean}{584}), which now carries $safety$. Since
\texttt{AddInduct} is \texttt{Type}-valued, that constructor is the one place where
$\mathtt{TrEnv}'$, itself still a \texttt{Prop}, quantifies over data that no consumer can read
back out.
\end{remark}

\begin{theorem}[A translated environment is well-formed]
\label{thm:tr-env-wf}
\contrib{mixed}
\lean{Lean4Lean.TrEnv'.wf}
\leanok
\uses{ind:tr-env}
\srcloc{Lean4Lean/Verify/Environment/Basic.lean}{591}
If $\mathtt{TrEnv}'\ safety\ C\ Q\ venv$ then $venv.\mathtt{WF}$: the model environment is built
by a legal sequence of \texttt{VEnv} declarations.
\end{theorem}
\begin{proof}
\uses{family:add-induct-stages, def:add-quot-witness}
\leanok
Induction on the relation; each case appends the corresponding \texttt{VEnv.Decls} step to the
sequence produced by the induction hypothesis. The \texttt{thm} case adds an \texttt{example}
step and then an axiom. The \texttt{quot} case uses \texttt{AddQuot.to\_addQuot}; the
\texttt{induct} case, contributed on \texttt{iota}, uses \texttt{AddInduct.env\_eq}.
\end{proof}

\section{Alignment and lookup transfer}

\begin{definition}[Alignment of the kernel map with the model's constants, \texttt{Aligned}]
\label{ind:aligned}
\contrib{mixed}
\lean{Lean4Lean.Aligned}
\leanok
\uses{def:tr-constant, def:insert-consts, def:venv-addpat}
\srcloc{Lean4Lean/Verify/Environment/Lemmas.lean}{12}
$\mathtt{Aligned}\ safety\ C\ venv$ is the weaker, constants-only companion of $\mathtt{TrEnv}'$:
it records that $C$ and $venv.\mathtt{constants}$ were built in step, by \texttt{empty},
\texttt{ignoreConst}, \texttt{const}, \texttt{defeq}, and the two \texttt{iota} additions
\texttt{pat} (registering a pattern reduction rule changes neither the map nor the constants) and
\texttt{block} (a whole list of fresh, pairwise distinct constants inserted at once, each
translating in the environment $venv'$ reached at the end of the block's \texttt{addConst} fold).
\thesisref{axioms.tex, \S ``Inductive types'' and \S ``The recursor''}
\end{definition}

\begin{remark}
The \texttt{block} constructor is the crux. A nested mutual block is inserted type by type by the
kernel, but the types of the block's recursors mention every type former and constructor of the
block, so it cannot be aligned one constant at a time. The constructor's docstring says exactly
that. \texttt{pat} is a one-line addition mirroring the new \texttt{VEnv.addPat}.
\end{remark}

\begin{lemma}[Basic consequences of alignment]
\label{family:aligned-basic}
\contrib{mixed}
\lean{Lean4Lean.Aligned.map_wf, Lean4Lean.Aligned.find?_iff, Lean4Lean.Aligned.find?, Lean4Lean.Aligned.find?_uniq}
\leanok
\srcloc{Lean4Lean/Verify/Environment/Lemmas.lean}{35}
\texttt{map\_wf}: the constant map is a well-formed \texttt{SMap}. \texttt{find?\_iff}: a name
resolves in $C$ to a declaration visible at $safety$ if and only if it resolves in
$venv.\mathtt{constants}$. \texttt{find?}: such a lookup produces a model constant that translates
the kernel one. \texttt{find?\_uniq}: a kernel lookup and a model lookup of the same name agree on
the name and are related by \texttt{TrConstant}.
\end{lemma}
\begin{proof}
\uses{ind:aligned, family:insert-consts-lemmas, family:indlem-addconst-spec}
\leanok
Induction on \texttt{Aligned}. The \texttt{pat} cases are monotonicity one-liners along
\texttt{VEnv.addPat\_le}; the \texttt{block} cases, which are substantial rather than
boilerplate, go through \texttt{insertConsts\_find?}, \texttt{VEnv.addConst\_foldlM\_find} and
\texttt{addConst\_foldlM\_constants\_inv}.
\end{proof}

\begin{lemma}[Quotient initialization preserves alignment]
\label{family:aligned-quot}
\lean{Lean4Lean.Aligned.addQuot1, Lean4Lean.Aligned.addQuot}
\leanok
\srcloc{Lean4Lean/Verify/Environment/Lemmas.lean}{75}
A continuation-passing combinator over one \texttt{AddQuot1} step, iterated four times, shows that
$\mathtt{AddQuot}\ C_1\ C_2\ venv_1\ venv_2$ carries $\mathtt{Aligned}\ safety\ C_1\ venv_1$ to
$\mathtt{Aligned}\ safety\ C_2\ venv_2$.
\end{lemma}
\begin{proof}
\uses{ind:aligned, def:add-quot-witness}
\leanok
Each step is an \texttt{Aligned.const} (at safety \texttt{safe}, lowered by \texttt{sf\_mono});
the final continuation adds the quotient defining equation by \texttt{Aligned.defeq}.
\end{proof}

\begin{lemma}[Registering a block's $\iota$ rules preserves alignment]
\label{thm:aligned-add-rules}
\contrib{iota}
\lean{Lean4Lean.Aligned.addRules}
\leanok
\srcloc{Lean4Lean/Verify/Environment/Lemmas.lean}{88}
If $decl.\mathtt{addRules}\ venv = \mathtt{some}\ venv'$ then alignment of $C$ with $venv$
transfers to $venv'$.
\end{lemma}
\begin{proof}
\uses{ind:aligned, family:indlem-foldlm, def:ind-add-rec-rule}
\leanok
Two nested \texttt{VEnv.foldlM\_inv} inductions over the recursors and their rules; each
\texttt{VEnv.addRecRule} step is unfolded to an \texttt{addPat} and closed by
\texttt{Aligned.pat}, which touches neither the map nor the constants.
\end{proof}

\begin{theorem}[Adding an inductive block preserves alignment]
\label{thm:aligned-add-induct}
\contrib{iota}
\lean{Lean4Lean.Aligned.addInduct}
\leanok
\uses{ind:aligned, struct:add-induct}
\srcloc{Lean4Lean/Verify/Environment/Lemmas.lean}{102}
From an $\mathtt{AddInduct}\ safety\ C_1\ venv_1\ decl\ C_2\ venv_2$ witness and
$\mathtt{Aligned}\ safety\ C_1\ venv_1$ one gets $\mathtt{Aligned}\ safety\ C_2\ venv_2$: the
three constant stages form a single \texttt{Aligned.block} step and the rule stage is
\ref{thm:aligned-add-rules}.
\thesisref{axioms.tex, \S ``Inductive types'' and \S ``The recursor''}
\end{theorem}
\begin{proof}
\uses{thm:aligned-add-rules, family:indlem-addconst-spec, thm:indlem-addtypesctorsrecs-eq, family:add-induct-names, family:add-induct-stages}
\leanok
Build the kind-wise matching of the kernel constants with \texttt{VInductDecl.consts}, all
translated in $env_R$ (monotonicity moves each from the stage where it was checked); permute it
into the kernel's insertion order with \texttt{List.Forall}$_2$\texttt{.perm\_left}; and show the
corresponding \texttt{addConst} fold succeeds by \texttt{VEnv.addConst\_foldlM\_perm} and
\texttt{VInductDecl.addTypesCtorsRecs\_eq}, an \texttt{addConst} fold being insensitive to order.
This replaces master's one-line \texttt{nomatch}, and isolates the whole kernel-order versus
model-order difficulty in one permutation argument.
\end{proof}

\begin{lemma}[Mutual blocks preserve alignment]
\label{family:aligned-defs}
\lean{Lean4Lean.Aligned.addDefEqs, Lean4Lean.Aligned.insertDefs}
\leanok
\srcloc{Lean4Lean/Verify/Environment/Lemmas.lean}{128}
Adding a list of defining equations leaves alignment untouched; inserting a block of definitions
whose headers translate keeps the map aligned with the extended environment.
\end{lemma}
\begin{proof}
\uses{ind:aligned, def:insert-defs}
\leanok
Both by induction on the block, the second threading freshness through the head insertion and
applying \texttt{Aligned.const} at each member.
\end{proof}

\begin{theorem}[A translated environment is aligned]
\label{thm:tr-env-aligned}
\contrib{mixed}
\lean{Lean4Lean.TrEnv'.aligned, Lean4Lean.TrEnv'.map_wf}
\leanok
\uses{ind:tr-env, ind:aligned}
\srcloc{Lean4Lean/Verify/Environment/Lemmas.lean}{162}
Every $\mathtt{TrEnv}'$ derivation yields an \texttt{Aligned} derivation for the same map and
environment; hence the kernel constant map of a translated environment is a well-formed
\texttt{SMap}.
\end{theorem}
\begin{proof}
\uses{thm:aligned-add-induct, family:aligned-quot, family:aligned-defs, family:aligned-basic}
\leanok
Induction on $\mathtt{TrEnv}'$, one constructor per case. The \texttt{induct} case, contributed on
\texttt{iota}, is \ref{thm:aligned-add-induct} applied to the induction hypothesis.
\end{proof}

\begin{lemma}[Pulling a lookup back across quotient initialization]
\label{family:quot-pullbacks}
\contrib{mixed}
\lean{Lean4Lean.pull_insert, Lean4Lean.AddQuot1.pull, Lean4Lean.AddQuot.pull}
\leanok
\srcloc{Lean4Lean/Verify/Environment/Lemmas.lean}{184}
\texttt{pull\_insert} pulls a lookup back across one insertion of a different constant.
\texttt{AddQuot1.pull} is the corresponding one-step combinator, and \texttt{AddQuot.pull}
iterates it four times: a constant that is not a \texttt{quotInfo} and is resolvable after
quotient initialization was already resolvable before.
\end{lemma}
\begin{proof}
\uses{def:add-quot-witness}
\leanok
\texttt{pull\_insert} rewrites by \texttt{SMap.find?\_insert} and discards the equal-name branch
by the assumed difference of the two constants; the iteration is the same combinator four times.
\end{proof}

\begin{remark}
Both branches contributed here, but not equally: \texttt{AddQuot1.pull} and \texttt{AddQuot.pull}
were introduced with the $\iota$ work and then generalised on \texttt{trproj}, which also added
\texttt{pull\_insert} for \ref{thm:tr-env-structure-rec}; per-line blame is 18 lines
\texttt{trproj} to 6 \texttt{iota}. The family is needed because
\ref{thm:tr-env-pats-iota} and \ref{thm:tr-env-structure-rec} induct over $\mathtt{TrEnv}'$ and
must survive the \texttt{quot} step.
\end{remark}

\begin{lemma}[A block of definitions preserves map well-formedness]
\label{thm:insert-defs-wf}
\lean{Lean4Lean.insertDefs_wf}
\leanok
\srcloc{Lean4Lean/Verify/Environment/Lemmas.lean}{211}
Inserting a list of \texttt{DefinitionVal}s with fresh, pairwise distinct names into a well-formed
map yields a well-formed map.
\end{lemma}
\begin{proof}
\uses{def:insert-defs}
\leanok
Induction on the block, reproving freshness of the tail after the head insertion. This is master
code relocated here from \texttt{Extension.lean} by the \texttt{iota} refactor, so the per-line
blame marks it as contributed.
\end{proof}

\begin{lemma}[\texttt{Aligned} lookup lemmas against the kernel environment]
\label{family:tr-env-find}
\lean{Lean4Lean.TrEnv.find?_iff, Lean4Lean.TrEnv.find?, Lean4Lean.TrEnv.find?_uniq}
\leanok
\srcloc{Lean4Lean/Verify/Environment/Lemmas.lean}{284}
The three \texttt{Aligned} lookup lemmas restated for $\mathtt{TrEnv}\ safety\ env\ venv$.
\end{lemma}
\begin{proof}
\uses{family:aligned-basic, thm:tr-env-aligned}
\leanok
Each converts \texttt{Environment.find?}, which is \texttt{SMap.find?'}, into \texttt{SMap.find?}
using \texttt{map\_wf}, and then applies the \texttt{Aligned} version.
\end{proof}

\begin{lemma}[Lookups in a mutual-definition block]
\label{family:insert-defs-lookup}
\lean{Lean4Lean.VEnv.addDefEqs_self, Lean4Lean.insertDefs_find?}
\leanok
\srcloc{Lean4Lean/Verify/Environment/Lemmas.lean}{304}
Every member of a block has its defining equation in the extended model environment; and a lookup
after \texttt{insertDefs} either predates the block or is one of the block's definitions.
\end{lemma}
\begin{proof}
\uses{def:insert-defs}
\leanok
Both by induction on the block.
\end{proof}

\begin{lemma}[Forward transport of a lookup across fresh insertions]
\label{family:forward-find-transport}
\contrib{iota}
\lean{Lean4Lean.insertDefs_find?_mono, Lean4Lean.AddQuot1.find?_mono, Lean4Lean.AddQuot.find?_mono}
\leanok
\srcloc{Lean4Lean/Verify/Environment/Lemmas.lean}{412}
A constant already resolvable stays resolvable to the same value across a block of definitions and
across quotient initialization.
\end{lemma}
\begin{proof}
\uses{def:insert-defs, def:add-quot-witness}
\leanok
\texttt{SMap.insertList\_find?\_mono} for the definition block, and
\texttt{SMap.find?\_insert\_of\_fresh} iterated four times through an \texttt{AddQuot1}
combinator. These are the forward duals of \ref{family:quot-pullbacks}; both directions are needed
because the induction hypotheses of \ref{thm:tr-env-pats-iota} and \ref{thm:tr-env-structure-rec}
carry lookups in both directions.
\end{proof}

\section{Soundness of delta reduction}

\begin{theorem}[$\delta$-reduction is sound in a translated environment]
\label{thm:tr-env-of-value}
\contrib{mixed}
\lean{Lean4Lean.TrEnv'.of_value, Lean4Lean.TrEnv.of_value}
\leanok
\uses{ind:tr-env}
\srcloc{Lean4Lean/Verify/Environment/Lemmas.lean}{333}
If a name resolves in a translated environment to a declaration $ci$ visible at $safety$ whose
$\delta$-value is $v$, then $v$ translates to something definitionally equal to the constant
itself: $\mathtt{TrExpr}\ venv\ ci.\mathtt{levelParams}\ []\ v\ (\mathtt{const}\ ci.\mathtt{name}\ \overline{\ell})$
with $\overline{\ell}$ the level parameters
of $ci$.
\thesisref{axioms.tex, \S ``Definitions ($\delta$ reduction)''}
\end{theorem}
\begin{proof}
\uses{family:add-induct-find, family:insert-defs-lookup}
\leanok
Induction on $\mathtt{TrEnv}'$. The \texttt{defn} and \texttt{mutualDef} cases apply
\texttt{VEnv.IsDefEq.extra0} to the registered defining equation; the \texttt{thm} case uses proof
irrelevance; the \texttt{axiom} and \texttt{opaque} cases are vacuous because those declarations
have no value; and the \texttt{induct} case, contributed on \texttt{iota}, is discharged by
\texttt{AddInduct.value\_find}, since the block's constants carry no $\delta$-value.
\end{proof}

\section{Structure-likeness}

\begin{lemma}[Model constants are never overwritten]
\label{thm:aligned-constants-pull}
\contrib{trproj}
\lean{Lean4Lean.Aligned.constants_pull}
\leanok
\srcloc{Lean4Lean/Verify/Environment/Lemmas.lean}{439}
If $C$ is aligned with $env$, $env \le env'$, and a name $c$ resolves in $C$ to a declaration
visible at $safety$, then a model constant for $c$ in $env'$ was already the model constant of
$env$.
\end{lemma}
\begin{proof}
\uses{family:aligned-basic}
\leanok
\texttt{find?\_iff} produces a model constant in $env$, monotonicity transports it into $env'$,
and injectivity of \texttt{some} identifies the two. Small but load-bearing: it is what lets
\ref{thm:tr-env-ctor-arity} read a constructor's arity off the environment stage where the
constructor was introduced rather than the final one.
\end{proof}

\begin{theorem}[A constructor's model type has arity $numParams + numFields$]
\label{thm:tr-env-ctor-arity}
\contrib{trproj}
\lean{Lean4Lean.TrEnv'.ctor_arity, Lean4Lean.TrEnv.ctor_arity}
\leanok
\uses{ind:tr-env, struct:tr-ind-type, def:pi-telescope}
\srcloc{Lean4Lean/Verify/Environment/Lemmas.lean}{449}
If a translated environment resolves $c$ to a kernel \texttt{ctorInfo} $cval$ and the model
resolves $c$ to $ci$, then
$ci.\mathtt{type}.\mathtt{piArity} = cval.\mathtt{numParams} + cval.\mathtt{numFields}$.
\thesisref{axioms.tex, \S ``Inductive types''}
\end{theorem}
\begin{proof}
\uses{thm:aligned-constants-pull, family:add-induct-find, family:quot-pullbacks, family:insert-defs-lookup}
\leanok
Visibility of $cval$ at $safety$ is recovered from the model lookup by
\texttt{Aligned.find?\_iff}. Then induction on $\mathtt{TrEnv}'$: every non-\texttt{induct} step
pulls the kernel lookup back through its insertion and the model lookup back with
\ref{thm:aligned-constants-pull}; the \texttt{induct} step, in the block case, reads the arity off
the \texttt{ctors} field of \texttt{TrIndType} after identifying the model constant through
\texttt{VEnv.addCtors\_find}.
\end{proof}

\begin{lemma}[\texttt{mkRecName} is injective]
\label{thm:mk-rec-name-inj}
\contrib{trproj}
\lean{Lean4Lean.mkRecName_inj}
\leanok
\srcloc{Lean4Lean/Verify/Environment/Lemmas.lean}{513}
$\mathtt{mkRecName}\ a = \mathtt{mkRecName}\ b$ implies $a = b$.
\end{lemma}
\begin{proof}
\leanok
\texttt{injection}: \texttt{mkRecName} appends the fixed component \texttt{rec}.
\end{proof}

\begin{theorem}[Structure-likeness: the recursor's telescope split]
\label{thm:tr-env-structure-rec}
\contrib{trproj}
\lean{Lean4Lean.TrEnv'.structure_rec, Lean4Lean.TrEnv.structure_rec}
\leanok
\uses{ind:tr-env, struct:tr-ind-type, struct:tr-recursor}
\srcloc{Lean4Lean/Verify/Environment/Lemmas.lean}{523}
Let a translated environment resolve $S$ to an \texttt{inductInfo} $ival$ and
$\mathtt{mkRecName}\ S$ to a \texttt{recInfo} $rval$, both visible at $safety$, and assume the
kernel's own structure tests $ival.\mathtt{all} = [S]$, $ival.\mathtt{ctors} = [ctorName]$ and
$ival.\mathtt{numIndices} = 0$. Then $rval.\mathtt{numParams} = ival.\mathtt{numParams}$,
$rval.\mathtt{numMotives} = 1$, $rval.\mathtt{numMinors} = 1$, $rval.\mathtt{numIndices} = 0$, and
$ctorName$ resolves to a \texttt{ctorInfo} $cval$ with
$cval.\mathtt{numParams} = ival.\mathtt{numParams}$.
\thesisref{axioms.tex, \S ``The recursor''}
\end{theorem}
\begin{proof}
\uses{thm:add-induct-ctor-find, family:quot-pullbacks, family:forward-find-transport, thm:mk-rec-name-inj, structure:ind-wf, family:add-induct-find, family:insert-defs-lookup}
\leanok
Induction on $\mathtt{TrEnv}'$, about 120 lines. Every non-\texttt{induct} step pulls both lookups back
and pushes the constructor lookup forward. The \texttt{induct} step splits four ways: if neither
constant comes from the block, apply the induction hypothesis; if exactly one does, the naming
field \texttt{TrRecursor.name\_major} together with \texttt{VInductDecl.WF.recs\_over\_block} or
\texttt{types\_have\_rec} produces a name that the freshness field of \texttt{AddInduct} says
cannot already resolve, a contradiction; if both do, the counts come from
\texttt{VInductDecl.WF.rec\_counts} and \texttt{rec\_params} together with the \texttt{all},
\texttt{numParams} and \texttt{numIndices} fields of \texttt{TrIndType}, and
\ref{thm:add-induct-ctor-find} supplies the constructor's parameter count.
\end{proof}

\reviewnote{Two caveats. The conclusion does not assert that the returned $cval$ is visible at
$safety$, so a caller that needs a \texttt{TrConstant} for it must re-derive visibility. And the
argument relies on the naming convention $\mathtt{mkRecName}\ S$ rather than on a structural link,
through the \texttt{TrRecursor.name\_major} field: a syntactic assumption about the kernel's
\texttt{AddInductive} that is asserted, not proved.}

\section{The iota-rule interface}

\begin{theorem}[Every kernel recursor rule is a registered $\iota$ rule]
\label{thm:tr-env-pats-iota}
\contrib{iota}
\lean{Lean4Lean.TrEnv'.pats_iota', Lean4Lean.TrEnv.pats_iota'}
\leanok
\uses{ind:tr-env, def:pattern-iota-rhs, ind:trexprs}
\srcloc{Lean4Lean/Verify/Environment/Lemmas.lean}{674}
If a translated environment resolves $recName$ to a \texttt{recInfo} $rval$ visible at $safety$
and $rval.\mathtt{rules}.\mathtt{find?}\ (\cdot.\mathtt{ctor} = cName)$ is $rule$, then there are a
constructor $cval$ resolvable under $cName$ and a closed model term $rhs$ translating
$rule.\mathtt{rhs}$ such that $venv.\mathtt{pats}$ contains the pattern
$\mathtt{SimplePattern.iota}\ recName\ M\ cName\ N$, with $M$ the kernel telescope split
$\mathtt{numParams} + \mathtt{numMotives} + \mathtt{numMinors} + \mathtt{numIndices}$ and
$N = cval.\mathtt{numParams} + rule.\mathtt{nfields}$, registered with reduct
$\mathtt{SimplePattern.iotaRHS}$ at those counts over $rhs$ and the trivial check.
\thesisref{axioms.tex, \S ``The computation rule ($\iota$ reduction)''}
\end{theorem}
\begin{proof}
\uses{thm:add-induct-rec-find, family:indlem-pat-registration, family:quot-pullbacks, family:forward-find-transport, family:add-induct-find}
\leanok
Induction on $\mathtt{TrEnv}'$. Every non-\texttt{induct} step transports the registration along
\texttt{VEnv.LE.pats} and the two lookups along the insertion lemmas, the \texttt{ignore} case
using the visibility hypothesis to rule out the inserted name. The \texttt{induct} step uses
\ref{thm:add-induct-rec-find} to resolve the recursor against the block and closes with
\texttt{VEnv.addInduct\_pat}. The \texttt{TrEnv} version additionally converts \texttt{SMap.find?}
into the environment's own \texttt{find?}.
\end{proof}

\reviewnote{The prime in \texttt{pats\_iota'} is a leftover of an earlier unprimed version that no
longer exists.}

\begin{definition}[The kernel data behind one registered $\iota$ pattern, \texttt{IotaRule}]
\label{struct:iota-rule}
\contrib{trproj}
\lean{Lean4Lean.TrEnv'.IotaRule}
\leanok
\uses{def:pattern-iota-rhs, def:ind-rec-shape, def:ind-rule-shape, def:ind-minor-for, ind:trexprs}
\srcloc{Lean4Lean/Verify/Environment/Lemmas.lean}{782}
$\mathtt{IotaRule}\ C\ venv\ recName\ cName\ M\ N\ r\ rval\ rule\ cval\ rhs\ hc\ rci$ bundles
eleven fields: $recName$ resolves in $C$ to \texttt{recInfo} $rval$, whose rule for $cName$ is
$rule$; $cName$ resolves to \texttt{ctorInfo} $cval$; $M$ is the kernel telescope split and
$N = cval.\mathtt{numParams} + rule.\mathtt{nfields}$; $rhs$ translates $rule.\mathtt{rhs}$; the
registered entry $r$ is \texttt{iotaRHS} at those counts, stated with \texttt{HEq} because the
type of $r.1$ mentions $M$ and $N$, with the trivial check; and on the model side $recName$ is a
registered constant $rci$ whose type has \texttt{RecShape} at that split and whose $j$-th minor
premise is a \texttt{MinorFor} $cName$ with at least $rule.\mathtt{nfields}$ binders, the reduct
having \texttt{RuleShape} at the matching counts.
\thesisref{axioms.tex, \S ``The computation rule ($\iota$ reduction)''}
\reviewnote{The \texttt{HEq} field is forced by the dependent \texttt{Pattern.RHS} type and is
honestly documented, but it makes the structure awkward to consume: the only consumer,
\ref{thm:tr-env-proj-defeq}, has to \texttt{cases eq\_of\_heq} it, and uses nine of the eleven
fields (all but \texttt{tr} and \texttt{rec\_shape}).}
\end{definition}

\begin{lemma}[Transport of an \texttt{IotaRule} along a non-\texttt{induct} step]
\label{thm:iota-rule-step}
\contrib{trproj}
\lean{Lean4Lean.TrEnv'.IotaRule.step}
\leanok
\srcloc{Lean4Lean/Verify/Environment/Lemmas.lean}{807}
Given $venv \le venv'$ and a map embedding that preserves lookups, an \texttt{IotaRule} over
$(C, venv)$ is one over $(C', venv')$: the kernel data is unchanged, only the translation and the
model constant lookup move.
\end{lemma}
\begin{proof}
\uses{struct:iota-rule, thm:trexprs-mono}
\leanok
Field by field, using \texttt{TrExprS.mono} for the translation and \texttt{VEnv.LE.constants}
for the recursor's model constant.
\end{proof}

\begin{theorem}[Inverse of \texttt{pats\_iota'}: a registered pattern comes from a kernel rule]
\label{thm:pats-iota-inv-shape}
\contrib{trproj}
\lean{Lean4Lean.TrEnv'.pats_iota_inv_shape}
\leanok
\uses{struct:iota-rule, ind:tr-env}
\srcloc{Lean4Lean/Verify/Environment/Lemmas.lean}{826}
If $venv.\mathtt{pats}$ contains
$(\mathtt{SimplePattern.iota}\ recName\ M\ cName\ N).\mathtt{toPattern}$ with entry $r$ in a
translated environment, then there exist $rval$, $rule$, $cval$, $rhs$, its closedness proof and
$rci$ with $\mathtt{IotaRule}\ C\ venv\ recName\ cName\ M\ N\ r\ rval\ rule\ cval\ rhs\ hc\ rci$.
In particular $M$ and $N$ are forced to be the kernel telescope split and
$cval.\mathtt{numParams} + rule.\mathtt{nfields}$.
\thesisref{axioms.tex, \S ``The computation rule ($\iota$ reduction)''}
\end{theorem}
\begin{proof}
\uses{thm:iota-rule-step, thm:add-induct-rec-reg, family:indlem-pats-origin, thm:indlem-iota-topattern-inj, family:indlem-addinduct-consts, family:add-induct-stages, family:quot-pullbacks, family:forward-find-transport, structure:ind-wf}
\leanok
Induction on $\mathtt{TrEnv}'$. Every non-\texttt{induct} step rewrites the \texttt{pats} field
away (\texttt{VEnv.addConst\_pats}, \texttt{addDefEq\_pats}, \texttt{addConsts\_pats},
\texttt{addQuot\_pats}) and applies \ref{thm:iota-rule-step}. The \texttt{induct} step uses
\texttt{VEnv.addInduct\_pats\_origin'} to decide whether the entry is old or new;
in the new case \texttt{VEnv.iota\_toPattern\_inj} recovers the key,
\ref{thm:add-induct-rec-reg} supplies the kernel side, and
\texttt{VInductDecl.WF.rec\_shape} and \texttt{rule\_shape} supply the model-side shape.
\end{proof}

\begin{remark}
This is the hardest theorem of the \texttt{trproj} branch in this group and the reason
\ref{thm:add-induct-rec-reg} exists. The statement is the right one:
\ref{thm:tr-env-proj-defeq} needs to know that a pattern it found through the \texttt{pat} field
of \texttt{TrProjCtor} really is the recursor's rule.
\end{remark}

\begin{lemma}[\texttt{pats\_iota\_inv\_shape} at the kernel environment level]
\label{thm:tr-env-pats-iota-inv-shape}
\contrib{trproj}
\lean{Lean4Lean.TrEnv.pats_iota_inv_shape}
\leanok
\srcloc{Lean4Lean/Verify/Environment/Lemmas.lean}{903}
The same statement with $C$ instantiated to $env.\mathtt{constants}$.
\end{lemma}
\begin{proof}
\uses{thm:pats-iota-inv-shape}
\leanok
Literally the application of \ref{thm:pats-iota-inv-shape}: the two statements are definitionally
equal.
\end{proof}

\reviewnote{Near-redundant. Unlike its sibling \ref{thm:tr-env-pats-iota}, it does not convert the
\texttt{SMap.find?} lookups of \texttt{IotaRule} into \texttt{Environment.find?}, so it adds
nothing over the primed version, and \ref{thm:tr-env-proj-defeq} still performs those conversions
by hand (\srcloc{Lean4Lean/Verify/Environment/Lemmas.lean}{1060} and
\srcloc{Lean4Lean/Verify/Environment/Lemmas.lean}{1072}). Its docstring claims the conversion it
does not perform.}

\begin{lemma}[A registered pattern fires on a well-typed redex]
\label{thm:tr-env-iota-defeq}
\contrib{iota}
\lean{Lean4Lean.TrEnv.iota_defeq}
\leanok
\srcloc{Lean4Lean/Verify/Environment/Lemmas.lean}{917}
If $venv.\mathtt{pats}\ p\ r$, a term $e$ matches $p$ with substitutions $m_1$, $m_2$, $e$ has type
$A$, and the check component of $r$ is realized by side conditions that are themselves
definitional equalities, then $e$ is definitionally equal (as an \texttt{IsDefEqU}) to the reduct
$r.1.\mathtt{apply}\ m_1\ m_2$.
\thesisref{axioms.tex, \S ``The computation rule ($\iota$ reduction)''}
\end{lemma}
\begin{proof}
\uses{rule:isdefeq-pat}
\leanok
Three lines: package \texttt{VEnv.IsDefEq.pat} at the type $A$ into \texttt{IsDefEqU}.
\end{proof}

\reviewnote{The name \texttt{TrEnv.iota\_defeq} places the lemma in the \texttt{TrEnv} namespace
although no \texttt{TrEnv} hypothesis appears in it; it is pure theory-layer material and belongs
next to \texttt{VEnv.IsDefEq.pat}.}

\begin{theorem}[The $\iota$ step of a translated environment]
\label{thm:tr-env-iota-rec}
\contrib{iota}
\lean{Lean4Lean.TrEnv.iota_rec}
\leanok
\uses{ind:tr-env, def:pattern-iota-rhs}
\srcloc{Lean4Lean/Verify/Environment/Lemmas.lean}{929}
If a translated environment resolves $recName$ to a visible \texttt{recInfo} $rval$ with a rule for
$cName$, resolves $cName$ to \texttt{ctorInfo} $cval$, and a term $e$ of type $A$ matches the
$\iota$ pattern at the kernel's telescope split and at
$cval.\mathtt{numParams} + rule.\mathtt{nfields}$, then there is a closed $rhs$ translating
$rule.\mathtt{rhs}$ with $e$ definitionally equal to the \texttt{iotaRHS} reduct applied to the
match.
\thesisref{axioms.tex, \S ``The computation rule ($\iota$ reduction)''}
\end{theorem}
\begin{proof}
\uses{thm:tr-env-pats-iota, thm:tr-env-iota-defeq}
\leanok
Compose \ref{thm:tr-env-pats-iota} with \ref{thm:tr-env-iota-defeq} at the trivial check, after
identifying the constructor returned by the former with the one given.
\end{proof}

\begin{remark}
This is the single point where the $\iota$ work of this group is consumed by the kernel model:
\texttt{inductiveReduceRecCore.WF} (declared at \srcloc{Lean4Lean/Verify/TypeChecker/WHNF.lean}{17},
calling \ref{thm:tr-env-iota-rec} at \srcloc{Lean4Lean/Verify/TypeChecker/WHNF.lean}{119}). That
theorem has no \texttt{sorry} of its own, though it inherits \texttt{sorryAx} through the strong
system. The enclosing \texttt{reduceRecursor.WF}
(\srcloc{Lean4Lean/Verify/TypeChecker/WHNF.lean}{147}, \texttt{sorry} at line 149) is not proved,
so the chain to \texttt{whnf} is not closed.
\end{remark}

\section{Beta telescopes and projection reduction}

\begin{lemma}[$\beta$-reduction of a saturated $\lambda$-telescope]
\label{family:beta-telescope}
\contrib{trproj}
\lean{Lean4Lean.VEnv.HasType.mkApps_inv_head, Lean4Lean.VEnv.IsDefEqU.beta_app, Lean4Lean.VEnv.IsDefEqU.mkApps_congr, Lean4Lean.VEnv.IsDefEqU.betaN}
\leanok
\srcloc{Lean4Lean/Verify/Environment/Lemmas.lean}{959}
For a well-formed environment and a well-typed application spine: the head of the spine is
well-typed (\texttt{mkApps\_inv\_head}); a $\lambda$ at the head $\beta$-reduces
(\texttt{beta\_app}); a definitional equality of heads extends along the spine
(\texttt{mkApps\_congr}); and, for a spine of exactly the telescope's arity,
$(\lambda x_1{:}A_1 \dots x_m{:}A_m.\,b)\ a_1 \dots a_m \equiv \mathtt{instFields}\ b\ [a_1,\dots,a_m]$
(\texttt{betaN}).
\thesisref{axioms.tex, \S ``Definitional equality'', the $\beta$ rule}
\axfoot{sorryAx via VEnv.WF.orderedStrong, VEnv.WF.patsStrong}
\end{lemma}
\begin{proof}
\uses{fam:strong-inversions, def:proj-inst-fields}
\leanok
Induction on the argument list. Each step inverts the application and the abstraction
(\texttt{HasType.app\_inv}, \texttt{HasType.lam\_inv}) to type the redex for
\texttt{IsDefEq.beta}, and uses unique typing (\texttt{IsDefEqU.of\_l}) to carry the head
equality along the spine; \texttt{betaN} then peels one binder at a time and closes by
transitivity.
\end{proof}

\reviewnote{This family is generic \texttt{VExpr} metatheory with no environment-translation
content; it sits in \texttt{Verify/Environment/Lemmas.lean} only because
\ref{thm:tr-env-proj-defeq} needs it, and belongs in the theory layer. It is also where the
projection work picks up \texttt{sorryAx}, but not, as one might expect, only through unique
typing: \texttt{mkApps\_inv\_head} uses nothing but \texttt{HasType.app\_inv}, whose
\texttt{OrderedStrong} argument is supplied by the \texttt{VEnv.WF} coercion
\texttt{VEnv.WF.orderedStrong}, and that is defined from the \texttt{sorry}
\texttt{VEnv.WF.patsStrong} (\srcloc{Lean4Lean/Theory/Typing/EnvLemmas.lean}{334}). Every lemma of
the family is tainted for that reason, whether or not it also uses \texttt{IsDefEqU.of\_l}.}

\begin{theorem}[Projection reduction]
\label{thm:tr-env-proj-defeq}
\contrib{trproj}
\lean{Lean4Lean.TrEnv.proj_defeq}
\leanok
\uses{ind:tr-env, struct:trprojctor, def:proj-field-selector}
\srcloc{Lean4Lean/Verify/Environment/Lemmas.lean}{1021}
Let a translated environment resolve $S$ to an \texttt{inductInfo} $ival$ with
$ival.\mathtt{all} = [S]$, $ival.\mathtt{ctors} = [ctorName]$ and $ival.\mathtt{numIndices} = 0$,
and resolve $ctorName$ to a \texttt{ctorInfo} $cval$. Let
$\mathtt{TrProjCtor}\ venv\ U\ \Gamma\ S\ i\ d\ e''\ ctorName\ usS\ uss\ params'\ np\ fieldTys$
present $e''$ as the recursor expansion of the $i$-th projection of $d$. If $\Gamma$ is a
well-formed context and $d$ is definitionally equal to a saturated constructor spine
$\mathtt{const}\ ctorName\ \overline{us}$ applied to $params \mathbin{+\!\!+} fields$ with
$|params| = ival.\mathtt{numParams}$, $|fields| = cval.\mathtt{numFields}$ and
$i < cval.\mathtt{numFields}$, then $e'' \equiv fields[i]$.
\thesisref{typesys.tex, \S \texttt{sec:undecidable}: the $\mathsf{inv}_x$ recursor idiom the
projection expansion copies}
\axfoot{sorryAx via VEnv.WF.orderedStrong, VEnv.WF.patsStrong; plus the three PersistentHashMap axioms}
\end{theorem}
\begin{proof}
\uses{thm:pats-iota-inv-shape, thm:tr-env-structure-rec, thm:tr-env-ctor-arity, family:beta-telescope, thm:pattern-iota-rhs-apply, thm:pattern-matches-varn-const, thm:proj-instfields-minor-spine}
\leanok
Six stages, labelled (A) to (F) in the source. (A) the constructor arity
$np + |fieldTys| = cval.\mathtt{numParams} + cval.\mathtt{numFields}$ from
\ref{thm:tr-env-ctor-arity} and \texttt{instPis\_piArity}. (B) the registered rule and its
model-side shape from \ref{thm:pats-iota-inv-shape}, after converting the two lookups by hand.
(C) the kernel's structure facts turned by \ref{thm:tr-env-structure-rec} into
$\mathtt{numMotives} = \mathtt{numMinors} = 1$ and $\mathtt{numIndices} = 0$. (D) the minor
premise has exactly the fields as binders, so the reduct has no recursive arguments. (E)
congruence from $d$ to the spine, then the $\iota$ step by \texttt{VEnv.IsDefEq.pat} with
\texttt{Pattern.matches\_varN\_const} on both sides and \texttt{iotaRHS'\_apply}. (F) two
\texttt{betaN} steps, one through the rule template and one through
\texttt{VExpr.fieldSelector}, landing on $fields[i]$ through
\texttt{instFields\_minor\_spine} and \texttt{instFields\_bvar}.
\end{proof}

\reviewnote{Technically the strongest result of this group, and correctly architected: the
projection computation rule is derived from the registered $\iota$ rule rather than postulated.
But it is orphaned. No non-test declaration uses it. The places that would use it are
\texttt{reduceProjCore.WF} (\srcloc{Lean4Lean/Verify/TypeChecker/Reduce.lean}{143}, \texttt{sorry}
at line 145; note that \texttt{reduceProj.WF} one line below it \emph{is} proved, by delegating to
\texttt{reduceProjCore.WF}), \texttt{inferProj.WF\_struct}
(\srcloc{Lean4Lean/Verify/TypeChecker/InferType.lean}{392}, \texttt{sorry} at line 398) and
\texttt{inferProj.WF} (\srcloc{Lean4Lean/Verify/TypeChecker/InferType.lean}{407}, \texttt{sorry} at
line 410). Its premise \texttt{TrProjCtor} is only ever inhabited by hand in
\texttt{Tests/ProjInhabit.lean}, the general construction \texttt{inferProj.WF\_struct} being
itself \texttt{sorry}. The parameter name
\texttt{kenv} also deviates from the \texttt{env} used by every neighbouring lemma.}

\begin{lemma}[Rebuilding an application spine in the translation]
\label{thm:tr-expr-mk-app-list}
\contrib{trproj}
\lean{Lean4Lean.TrExpr.mkAppList}
\leanok
\srcloc{Lean4Lean/Verify/Environment/Lemmas.lean}{1166}
If a head and a list of arguments translate, and the model spine $f'.\mathtt{mkApps}\ as'$ is
well-typed, then $f.\mathtt{mkAppList}\ as$ translates to $f'.\mathtt{mkApps}\ as'$.
\axfoot{sorryAx via VEnv.WF.orderedStrong, VEnv.WF.patsStrong}
\end{lemma}
\begin{proof}
\uses{family:beta-telescope, family:trexpr-builders}
\leanok
Induction on the pointwise translation of the arguments; each application node needs the typing of
its function and argument, which \texttt{HasType.mkApps\_inv\_head} and \texttt{HasType.app\_inv}
read off the whole spine. Unlike \ref{thm:tr-env-proj-defeq} this lemma does reach the kernel
model, at \srcloc{Lean4Lean/Verify/TypeChecker/WHNF.lean}{131}.
\end{proof}

\section{Quotient initialization}

\begin{lemma}[Unfolding one \texttt{withLocalDecl} of the expression builder]
\label{thm:with-local-decl-run}
\contrib{iota}
\lean{Lean4Lean.withLocalDecl_run}
\leanok
\srcloc{Lean4Lean/Verify/Environment/Quot.lean}{20}
$\mathtt{withLocalDecl}\ n\ bi\ ty\ k$, run at a local context and a name generator, equals $k$
applied to the fresh free variable $\langle ngen.\mathtt{curr}\rangle$, the context extended by
\texttt{mkLocalDecl}, and $ngen.\mathtt{next}$.
\end{lemma}
\begin{proof}
\leanok
\texttt{rfl}.
\end{proof}

\reviewnote{Used once, inside \ref{thm:check-eq-type-ok}, but declared outside the
\texttt{AddQuotAux} namespace, so a very general name lands in the \texttt{Lean4Lean} root
namespace.}

\begin{definition}[A symbolic replay of what \texttt{Environment.addQuot} builds]
\label{def:quot-telescope-data}
\contrib{iota}
\lean{Lean4Lean.AddQuotAux.ng0, Lean4Lean.AddQuotAux.x1, Lean4Lean.AddQuotAux.x2, Lean4Lean.AddQuotAux.x3, Lean4Lean.AddQuotAux.x4, Lean4Lean.AddQuotAux.x5, Lean4Lean.AddQuotAux.x6, Lean4Lean.AddQuotAux.u, Lean4Lean.AddQuotAux.v, Lean4Lean.AddQuotAux.αr, Lean4Lean.AddQuotAux.quot_r, Lean4Lean.AddQuotAux.L1, Lean4Lean.AddQuotAux.L2, Lean4Lean.AddQuotAux.L3, Lean4Lean.AddQuotAux.L2', Lean4Lean.AddQuotAux.L3', Lean4Lean.AddQuotAux.L4, Lean4Lean.AddQuotAux.L5, Lean4Lean.AddQuotAux.L6, Lean4Lean.AddQuotAux.L4i, Lean4Lean.AddQuotAux.L5i, Lean4Lean.AddQuotAux.T1, Lean4Lean.AddQuotAux.T2, Lean4Lean.AddQuotAux.sanity, Lean4Lean.AddQuotAux.T3, Lean4Lean.AddQuotAux.all_quot, Lean4Lean.AddQuotAux.T4, Lean4Lean.AddQuotAux.q1, Lean4Lean.AddQuotAux.q2, Lean4Lean.AddQuotAux.q3, Lean4Lean.AddQuotAux.q4}
\leanok
\srcloc{Lean4Lean/Verify/Environment/Quot.lean}{27}
Abbreviations reproducing, symbolically, the state the kernel's quotient initializer reaches: the
six free variables its name generator produces, the levels $u$ and $v$, the abbreviations
$\alpha r$ for the relation type and \texttt{quot\_r} for $\mathtt{Quot}\ \alpha\ r$, the ten local
contexts
$L_1, \dots, L_{5i}$ built by \texttt{mkLocalDecl}, the four \texttt{mkForall} telescopes
$T_1, \dots, T_4$ (with the auxiliary bodies \texttt{sanity} for \texttt{Quot.lift} and
\texttt{all\_quot} for \texttt{Quot.ind}), and the four resulting \texttt{quotInfo} declarations
$q_1, \dots, q_4$ for \texttt{Quot}, \texttt{Quot.mk}, \texttt{Quot.lift}, \texttt{Quot.ind}.
\thesisref{axioms.tex, \S ``Quotient types''}
\reviewnote{These abbreviations mirror \texttt{Lean4Lean/Quot.lean} line by line, which is the
only way to compute with a monadic expression builder, but nothing links the two beyond the
proofs breaking: any change to a binder name, a binder info or the order of the
\texttt{withLocalDecl} calls silently invalidates \ref{thm:add-quot-eq}.}
\end{definition}

\begin{lemma}[Lookup ladder for the quotient telescope contexts]
\label{family:quot-telescope-find}
\contrib{iota}
\lean{Lean4Lean.AddQuotAux.L1_mwf, Lean4Lean.AddQuotAux.L1_find, Lean4Lean.AddQuotAux.L2_mwf, Lean4Lean.AddQuotAux.L2_find, Lean4Lean.AddQuotAux.L3_find, Lean4Lean.AddQuotAux.L2'_mwf, Lean4Lean.AddQuotAux.L2'_find, Lean4Lean.AddQuotAux.L3'_mwf, Lean4Lean.AddQuotAux.L3'_find, Lean4Lean.AddQuotAux.L4_mwf, Lean4Lean.AddQuotAux.L4_find, Lean4Lean.AddQuotAux.L5_mwf, Lean4Lean.AddQuotAux.L5_find, Lean4Lean.AddQuotAux.L6_find, Lean4Lean.AddQuotAux.L4i_mwf, Lean4Lean.AddQuotAux.L4i_find, Lean4Lean.AddQuotAux.L5i_find}
\leanok
\srcloc{Lean4Lean/Verify/Environment/Quot.lean}{92}
For each context $L$: the underlying \texttt{fvarIdToDecl} map is well-formed, and
$L.\mathtt{find?}\ x$ reduces to one decidable comparison against the newest free variable
followed by the previous context's lookup.
\end{lemma}
\begin{proof}
\uses{def:quot-telescope-data, family:lctx-empty-decl}
\leanok
All by \texttt{LocalContext.find?\_mkLocalDecl}, which needs only well-formedness of the
underlying map, threaded along the chain from the empty context.
\end{proof}

\begin{definition}[Two macros driving the telescope computations]
\label{def:quot-tactics}
\contrib{iota}
\lean{Lean4Lean.AddQuotAux.tacticQuot_simp, Lean4Lean.AddQuotAux.tacticQuot_mem}
\leanok
\uses{family:quot-telescope-find}
\srcloc{Lean4Lean/Verify/Environment/Quot.lean}{144}
\texttt{quot\_simp} rewrites along the lookup ladder, decides the free-variable comparisons and
unfolds \texttt{mkBindingList1} and \texttt{Expr.abstract}. \texttt{quot\_mem} discharges the side
condition of \texttt{LocalContext.mkForall\_eq\_fold} that every telescope variable is declared in
its context. The two are \texttt{macro} declarations; their elaborated constants carry the
\texttt{tactic} prefix.
\end{definition}

\begin{lemma}[Evaluating the four quotient types to closed expressions]
\label{family:quot-type-eqs}
\contrib{iota}
\lean{Lean4Lean.AddQuotAux.T1_eq, Lean4Lean.AddQuotAux.T2_eq, Lean4Lean.AddQuotAux.sanity_eq, Lean4Lean.AddQuotAux.T3_eq, Lean4Lean.AddQuotAux.all_quot_eq, Lean4Lean.AddQuotAux.T4_inner_eq, Lean4Lean.AddQuotAux.T4_eq}
\leanok
\srcloc{Lean4Lean/Verify/Environment/Quot.lean}{159}
Each \texttt{mkForall} telescope is computed to an explicit de Bruijn \texttt{Expr}. For instance
$T_1 = \forall \{\alpha : \mathsf{Sort}\ u\}\ (r : \alpha \to \alpha \to \mathsf{Prop}),\ \mathsf{Sort}\ u$,
and analogously for \texttt{Quot.mk}, for \texttt{Quot.lift} through the
\texttt{sanity} hypothesis, and for \texttt{Quot.ind} through \texttt{all\_quot}.
\thesisref{axioms.tex, \S ``Quotient types''}
\end{lemma}
\begin{proof}
\uses{def:quot-telescope-data, def:quot-tactics, thm:mkforall-eq-fold}
\leanok
Rewrite \texttt{LocalContext.mkForall} into a fold by \texttt{mkForall\_eq\_fold}, whose side
conditions are discharged by \texttt{quot\_mem}, and evaluate the fold by \texttt{quot\_simp}.
\end{proof}

\begin{lemma}[The computed effect of \texttt{Environment.addQuot}]
\label{thm:add-quot-eq}
\contrib{iota}
\lean{Lean4Lean.AddQuotAux.addQuot_eq}
\leanok
\srcloc{Lean4Lean/Verify/Environment/Quot.lean}{72}
If quotient initialization has not happened, \texttt{checkEqType} succeeds, and the four names
pass \texttt{checkName}, then $\mathtt{Environment.addQuot}\ env$ returns
\texttt{markQuotInit} of $env$ with $q_1, q_2, q_3, q_4$ added in order.
\thesisref{axioms.tex, \S ``Quotient types''}
\end{lemma}
\begin{proof}
\uses{def:quot-telescope-data}
\leanok
Unfold \texttt{Environment.addQuot}, rewrite by the five hypotheses, and close by \texttt{rfl}.
\end{proof}

\begin{definition}[The expected type of \texttt{Eq}]
\label{def:check-eq-type-shape}
\contrib{iota}
\lean{Lean4Lean.AddQuotAux.LE1, Lean4Lean.AddQuotAux.TE, Lean4Lean.AddQuotAux.LE1_find, Lean4Lean.AddQuotAux.TE_eq}
\leanok
\uses{def:quot-tactics, thm:mkforall-eq-fold}
\srcloc{Lean4Lean/Verify/Environment/Quot.lean}{224}
$\mathtt{TE}\ u$ is the telescope
$\forall \{\alpha : \mathsf{Sort}\ u\},\ \alpha \to \alpha \to \mathsf{Prop}$ exactly as
\texttt{checkEqType} builds it, \texttt{LE1\_find} is its one-entry lookup lemma, and
\texttt{TE\_eq} evaluates it to the explicit de Bruijn expression.
\thesisref{axioms.tex, \S ``Quotient types''}
\end{definition}

\begin{lemma}[\texttt{Environment.get} success implies \texttt{find?} success]
\label{thm:environment-get-ok}
\contrib{iota}
\lean{Lean4Lean.AddQuotAux.Environment.get_ok}
\leanok
\srcloc{Lean4Lean/Verify/Environment/Quot.lean}{244}
$env.\mathtt{get}\ n = \mathtt{ok}\ ci$ implies $env.\mathtt{find?}\ n = \mathtt{some}\ ci$.
\end{lemma}
\begin{proof}
\leanok
Case split on \texttt{find?}: the \texttt{none} branch contradicts the success of \texttt{get}.
\end{proof}

\reviewnote{Declared inside \texttt{AddQuotAux}, so this generally useful lemma has the full name
\texttt{Lean4Lean.AddQuotAux.Environment.get\_ok} and is hidden in an auxiliary namespace.}

\begin{theorem}[What a successful \texttt{checkEqType} establishes]
\label{thm:check-eq-type-ok}
\contrib{iota}
\lean{Lean4Lean.AddQuotAux.checkEqType_ok}
\leanok
\uses{def:check-eq-type-shape}
\srcloc{Lean4Lean/Verify/Environment/Quot.lean}{253}
If $\mathtt{checkEqType}\ env = \mathtt{ok}\ ()$ then \texttt{Eq} resolves in $env$ to an
\texttt{inductInfo} that is not unsafe, has exactly one universe parameter $u$, and whose type is
syntactically $\mathtt{TE}\ u$.
\thesisref{axioms.tex, \S ``Quotient types''}
\end{theorem}
\begin{proof}
\uses{thm:environment-get-ok, thm:with-local-decl-run}
\leanok
Destruct the monadic check: the constant must be an \texttt{inductInfo}, not unsafe, with a
one-element \texttt{levelParams} and a one-element \texttt{ctors} list; the type equality is
extracted from the \texttt{bne} branch after unfolding \texttt{ExprBuildT} and
\ref{thm:with-local-decl-run}.
\end{proof}

\begin{remark}
This replaces master's \texttt{checkEqType.WF}, which proved \texttt{False}: on master no
\texttt{TrEnv} could contain the inductive \texttt{Eq}, because \texttt{AddInduct} had no
constructors, so quotient initialization was vacuously well-formed in the interesting branch.
Making \texttt{AddInduct} inhabitable forced this file into existence, and it is the clearest
example of the \texttt{iota} contribution paying for itself.
\end{remark}

\begin{lemma}[The quotient types translate to the model's]
\label{family:quot-type-tr}
\contrib{iota}
\lean{Lean4Lean.AddQuotAux.T1', Lean4Lean.AddQuotAux.T2', Lean4Lean.AddQuotAux.T3', Lean4Lean.AddQuotAux.T4', Lean4Lean.AddQuotAux.T1_tr, Lean4Lean.AddQuotAux.T2_tr, Lean4Lean.AddQuotAux.T3_tr, Lean4Lean.AddQuotAux.T4_tr}
\leanok
\srcloc{Lean4Lean/Verify/Environment/Quot.lean}{305}
Each computed kernel type $T_i'$ has a \texttt{TrExprS} derivation into the type of the
corresponding model constant (\texttt{quotConst}, \texttt{quotMkConst}, \texttt{quotLiftConst},
\texttt{quotIndConst}), given that the constants it mentions (\texttt{Quot}, \texttt{Quot.mk},
\texttt{Eq}) are already registered in the model.
\thesisref{axioms.tex, \S ``Quotient types''}
\end{lemma}
\begin{proof}
\uses{ind:trexprs, def:quot-consts, family:quot-type-eqs}
\leanok
Explicit \texttt{TrExprS} derivation trees, built binder by binder with \texttt{refine .forallE},
every typing side condition discharged by the \texttt{type\_tac} macro.
\end{proof}

\reviewnote{Sixty lines of hand-built derivation trees: correct, but unreadable and fragile. A
\texttt{TrExprS}-building tactic would collapse them, as the \texttt{type\_tac} macro they rely on
already notes for its own case.}

\begin{theorem}[A well-shaped \texttt{Eq} is the model's \texttt{eqConst}]
\label{thm:quot-ready}
\contrib{iota}
\lean{Lean4Lean.AddQuotAux.quotReady}
\leanok
\uses{ind:tr-env, def:venv-addquot}
\srcloc{Lean4Lean/Verify/Environment/Quot.lean}{398}
If a translated environment resolves \texttt{Eq} to a safe \texttt{inductInfo} with one level
parameter whose type is $\mathtt{TE}\ u$, then $venv.\mathtt{QuotReady}$: the model environment's
\texttt{Eq} is exactly \texttt{eqConst}.
\thesisref{axioms.tex, \S ``Quotient types''}
\end{theorem}
\begin{proof}
\uses{thm:check-eq-type-ok, family:tr-env-find, def:check-eq-type-shape}
\leanok
\texttt{TrEnv.find?} yields a model constant translating \texttt{Eq}; \texttt{TrExprS.eqv}
transports the derivation along the syntactic type equality, the derivation is then inverted
\texttt{forallE}, \texttt{sort}, \texttt{bvar} by \texttt{cases}, and the level and de Bruijn side
conditions pin the model type to \texttt{eqConst}'s.
\end{proof}

\begin{theorem}[The \texttt{AddQuot} witness of quotient initialization]
\label{thm:exists-add-quot}
\contrib{iota}
\lean{Lean4Lean.AddQuotAux.C', Lean4Lean.AddQuotAux.exists_addQuot}
\leanok
\uses{def:add-quot-witness, ind:tr-env}
\srcloc{Lean4Lean/Verify/Environment/Quot.lean}{436}
If a translated environment is \texttt{QuotReady} and the four quotient names are fresh, then the
four model \texttt{addConst}s succeed in sequence and the resulting environment, extended by
\texttt{quotDefEq}, is related to the kernel map $C'$ (the four insertions) by \texttt{AddQuot}.
\thesisref{axioms.tex, \S ``Quotient types''}
\end{theorem}
\begin{proof}
\uses{family:quot-type-eqs, family:quot-type-tr, thm:quot-ready, def:quot-consts}
\leanok
Construct the four intermediate environments by \texttt{VEnv.addConst\_eq\_none}, propagate
freshness through \texttt{addConst\_eq\_of\_ne}, and supply each \texttt{AddQuot1} step with its
\texttt{TrConstant}, obtained from the computed type equalities and the translations of
\ref{family:quot-type-tr}.
\end{proof}

\begin{theorem}[Quotient initialization preserves well-formedness]
\label{thm:add-quot-wf}
\contrib{iota}
\lean{Lean4Lean.addQuot.WF}
\leanok
\uses{ind:tr-env, def:vtc-venvs}
\srcloc{Lean4Lean/Verify/Environment/Quot.lean}{500}
A successful $\mathtt{Environment.addQuot}\ env$ yields a model family $ves'$ that is well-formed
for the new environment and extends $ves$ at every safety level. If the environment was already
initialized nothing happens; otherwise the four quotient constants and the computation rule are
added by the \texttt{quot} step of $\mathtt{TrEnv}'$.
\thesisref{axioms.tex, \S ``Quotient types''}
\end{theorem}
\begin{proof}
\uses{thm:exists-add-quot, thm:quot-ready, thm:add-quot-eq, family:check-name, family:safe-primitives}
\leanok
Extract the five successful checks from the monadic result, rewrite by \ref{thm:add-quot-eq},
turn the shape of \texttt{Eq} into \texttt{VEnv.QuotReady} at every level by
\ref{thm:quot-ready}, choose the extended model family with \texttt{VEnvs.axiom\_of\_choice}, and
rebuild the four \texttt{VEnvs.WF} fields (\texttt{tr}, \texttt{hasPrimitives},
\texttt{safePrimitives}, \texttt{mono}) one by one.
\end{proof}

\reviewnote{Two smells in the assembly: the four \texttt{checkName} results are extracted by four
near-identical seven-line blocks
(\srcloc{Lean4Lean/Verify/Environment/Quot.lean}{521}), and the \texttt{safePrimitives} field is
rebuilt by a 22-line hand-nested chain of \texttt{Environment.find?\_add\_of\_ne}
(\srcloc{Lean4Lean/Verify/Environment/Quot.lean}{584}); a single auxiliary lemma would remove
both. Note also that \texttt{Quot.sound}, the only genuine axiom among the quotient constants in
the thesis, is not added here (Lean declares it in core), so this theorem covers the four
kernel-builtin constants plus the computation rule.}

\section{The extension machinery}

\begin{lemma}[Monotonicity of model extension, and existence of extensions]
\label{family:venv-mono}
\contrib{mixed}
\lean{Lean4Lean.VEnv.addConst_mono, Lean4Lean.VEnv.addDefEq_mono, Lean4Lean.VEnv.addConsts_mono, Lean4Lean.VEnv.addDefEqs_mono, Lean4Lean.TrEnv.exists_addConst, Lean4Lean.TrEnv.constants_eq_none, Lean4Lean.TrEnv.exists_addConsts}
\leanok
\srcloc{Lean4Lean/Verify/Environment/Extension.lean}{11}
Adding the same constant, defining equation, or list of either, to two environments related by
$\le$ preserves $\le$; and in a translated environment a kernel name that does not resolve has no
model constant, so a fresh name can always be given one.
\end{lemma}
\begin{proof}
\uses{def:venv-add-const, family:tr-env-find, def:venv-basic}
\leanok
By unfolding \texttt{addConst} and its list form, using \texttt{find?\_iff} to rule out a clash.
The \texttt{iota} branch touched only two lines here, adding the \texttt{pats} component to two
\texttt{VEnv.LE} constructions when \texttt{VEnv.LE} gained that field.
\end{proof}

\begin{lemma}[Bookkeeping for the primitive-declaration side conditions]
\label{family:safe-primitives}
\lean{Lean4Lean.safePrimitives_add', Lean4Lean.VEnvs.WF.safePrimitives_add, Lean4Lean.VEnvAt.safePrimitives_add, Lean4Lean.VEnv.HasPrimitives.addConsts, Lean4Lean.VEnv.HasPrimitives.addDefEqs, Lean4Lean.Environment.constants_addDefs, Lean4Lean.VEnvs.WF.safePrimitives_addDefs, Lean4Lean.Environment.quotInit_addDefs, Lean4Lean.Environment.find?_add_of_ne}
\leanok
\srcloc{Lean4Lean/Verify/Environment/Extension.lean}{46}
The \texttt{safePrimitives} invariant (every name in \texttt{Environment.primitives} is declared
safe and with no level parameters) and the \texttt{HasPrimitives} model invariant are preserved by
each way of extending the environment, together with the constant-map bookkeeping for definition
blocks.
\end{lemma}
\begin{proof}
\leanok
Each is a short computation on the constant map or on the model environment, keyed on the name
being distinct from the primitive in question.
\end{proof}

\begin{lemma}[A block invisible at a safety level extends only the map]
\label{thm:tr-env-ignore-defs}
\lean{Lean4Lean.TrEnv'.ignoreDefs}
\leanok
\srcloc{Lean4Lean/Verify/Environment/Extension.lean}{135}
A list of definitions none of which is visible at $safety$ extends the constant map by
\texttt{insertDefs} without changing the model environment.
\end{lemma}
\begin{proof}
\uses{ind:tr-env, def:insert-defs}
\leanok
Induction on the list, applying the \texttt{ignore} step of $\mathtt{TrEnv}'$ at each member.
\end{proof}

\begin{definition}[Headers of a mutual block, and the temporary environment]
\label{def:tr-mutual-header}
\lean{Lean4Lean.TrMutualHeader, Lean4Lean.VEnvAt.addAxioms}
\leanok
\uses{def:tr-constant}
\srcloc{Lean4Lean/Verify/Environment/Extension.lean}{159}
\texttt{TrMutualHeader} is the data the header loop of \texttt{addMutual} produces per member: a
translated \texttt{VDefVal} header, its well-formedness, freshness of the name, and
non-primitiveness. \texttt{VEnvAt.addAxioms} models the temporary environment in which the bodies
are checked, in which every header has been added as an axiom, so a body may refer to any member
of the block, itself included, but cannot $\delta$-unfold it.
\thesisref{axioms.tex, \S ``Definitions ($\delta$ reduction)''}
\end{definition}

\begin{lemma}[Adding a whole mutual block preserves well-formedness]
\label{thm:add-mutual-block-wf}
\lean{Lean4Lean.addMutualBlock.WF}
\leanok
\srcloc{Lean4Lean/Verify/Environment/Extension.lean}{209}
Given headers checked in $env$, bodies checked in the temporary environment, distinct fresh names
and non-primitiveness, the environment obtained by adding all the definitions is well-formed and
extends the old one at every safety level.
\end{lemma}
\begin{proof}
\uses{ind:tr-env, def:tr-mutual-header, thm:tr-env-ignore-defs, family:safe-primitives}
\leanok
Split on visibility of the block at each safety level: where it is invisible,
\ref{thm:tr-env-ignore-defs}; where it is visible, the \texttt{mutualDef} step of
$\mathtt{TrEnv}'$, with the primitive side conditions discharged by \ref{family:safe-primitives}.
Like the unsafe branch of \texttt{addDefinition}, this cannot conclude a \texttt{VEnv.AddDef} for
the individual members, because their bodies do not translate before the block is added.
\end{proof}

\begin{lemma}[Generic single-constant extension lemmas]
\label{family:add-const-wf}
\lean{Lean4Lean.addConstCore.WF, Lean4Lean.addConst.WF, Lean4Lean.addDef.WF, Lean4Lean.addUnsafeDef.WF}
\leanok
\srcloc{Lean4Lean/Verify/Environment/Extension.lean}{291}
The reusable core behind \ref{thm:add-axiom-wf}, \ref{thm:add-theorem-wf},
\ref{thm:add-opaque-wf} and \ref{thm:add-definition-wf}: given a translation of the header,
well-formedness of the model constant, freshness, non-primitiveness and a per-safety-level
$\mathtt{TrEnv}'$ step, adding the constant, and optionally its defining equation, preserves
\texttt{VEnvs.WF} and produces the corresponding \texttt{VEnv.AddConst} or \texttt{VEnv.AddDef}
step.
\end{lemma}
\begin{proof}
\uses{def:venv-add-const, ind:tr-env, family:safe-primitives, def:vtc-venvs}
\leanok
\texttt{VEnvs.axiom\_of\_choice} picks the extended family, then the four \texttt{VEnvs.WF} fields
are rebuilt one by one.
\end{proof}

\section{Per-declaration checks}

\begin{lemma}[Name freshness and metavariable or free-variable checks]
\label{family:check-name}
\lean{Lean4Lean.ConstantInfo.defnInfo_safety, Lean4Lean.checkName.WF, Lean4Lean.checkNoMVar.WF, Lean4Lean.checkNoFVar.WF, Lean4Lean.checkNoMVarNoFVar.WF, Lean4Lean.Except.WF.trivial, Lean4Lean.TypeChecker.M.WF.pureBind}
\leanok
\srcloc{Lean4Lean/Verify/Environment/Checker.lean}{8}
A successful \texttt{checkName} means the name does not resolve in the environment and, unless
explicitly allowed, is not a primitive. A successful \texttt{checkNoMVarNoFVar} means the
expression has no metavariables and no free variables.
\end{lemma}
\begin{proof}
\leanok
Direct inspection of the monadic checks, through the \texttt{Except.WF} and
\texttt{TypeChecker.M.WF} framework.
\end{proof}

\begin{lemma}[Verification of the shared declaration-header check]
\label{family:check-constant-val}
\lean{Lean4Lean.checkConstantValBody.WF, Lean4Lean.checkConstantValCore.WF, Lean4Lean.checkConstantVal.WF}
\leanok
\srcloc{Lean4Lean/Verify/Environment/Checker.lean}{56}
A successful \texttt{checkConstantVal} on a \texttt{ConstantInfo} produces a model
\texttt{VConstVal} translating the header, its well-formedness, freshness of the name and
non-primitiveness: the common precondition of every \texttt{add} lemma.
\end{lemma}
\begin{proof}
\uses{def:tr-constant, family:check-name}
\leanok
Runs the level-parameter duplication check, \texttt{checkNoMVarNoFVar} and \texttt{checkType}
through the \texttt{TypeChecker.M.WF} framework; the translation comes from
\texttt{checkType.WF}, whose correctness is where this family inherits \texttt{sorryAx}.
\end{proof}

\begin{lemma}[Verification of the body check]
\label{family:check-body}
\lean{Lean4Lean.checkBodyCore, Lean4Lean.checkBodyCore.WF, Lean4Lean.checkBody, Lean4Lean.checkBody.WF}
\leanok
\srcloc{Lean4Lean/Verify/Environment/Checker.lean}{104}
A successful body check yields a model term translating the declaration's value together with a
\texttt{HasType} derivation against the declared type.
\end{lemma}
\begin{proof}
\uses{family:check-constant-val}
\leanok
\texttt{TypeChecker.checkType.WF} translates the value, then \texttt{isDefEq.WF} converts its
inferred type to the declared one.
\end{proof}

\begin{lemma}[Per-declaration-kind checks and their verification]
\label{family:check-decls}
\lean{Lean4Lean.checkTheorem, Lean4Lean.checkTheorem.WF, Lean4Lean.checkDefinitionBody.WF, Lean4Lean.checkDefinition, Lean4Lean.checkDefinition.WF, Lean4Lean.checkOpaque, Lean4Lean.checkOpaque.WF}
\leanok
\srcloc{Lean4Lean/Verify/Environment/Checker.lean}{146}
Reimplementations of the kernel's theorem, definition and opaque checks in the verified
\texttt{TypeChecker} monad, each with a \texttt{WF} lemma producing exactly the data its
consumer needs: header translation, body translation, Prop-ness for theorems, and
primitive preservation for definitions.
\end{lemma}
\begin{proof}
\uses{family:check-constant-val, family:check-body}
\leanok
Each composes \ref{family:check-constant-val} with \ref{family:check-body} and, for theorems, the
\texttt{Sort 0} check on the statement.
\end{proof}

\section{The top-level driver}

\begin{lemma}[Adding an axiom preserves well-formedness]
\label{thm:add-axiom-wf}
\lean{Lean4Lean.addAxiom.WF}
\leanok
\srcloc{Lean4Lean/Verify/Environment.lean}{11}
If the safety-indexed model family $ves$ is well-formed for $env$, then a successful
$\mathtt{addAxiom}\ env\ v$ yields an environment for which there is a well-formed family $ves'$
and a model constant $ci'$ such that at every safety level the step is a \texttt{VEnv.AddConst} of
$\mathtt{axiomInfo}\ v$.
\thesisref{axioms.tex, \S ``Definitions ($\delta$ reduction)''}
\end{lemma}
\begin{proof}
\uses{family:check-constant-val, family:add-const-wf, def:venv-add-const, ind:tr-env}
\leanok
\texttt{checkConstantVal.WF} translates the header at the safety level the kernel checks at, then
\texttt{addConst.WF} with the \texttt{axiom} step of $\mathtt{TrEnv}'$.
\end{proof}

\begin{lemma}[Adding a definition preserves well-formedness]
\label{thm:add-definition-wf}
\lean{Lean4Lean.addDefinition.WF}
\leanok
\srcloc{Lean4Lean/Verify/Environment.lean}{28}
A successful $\mathtt{addDefinition}\ env\ v$ yields a well-formed model family extending the old
one at every safety level, and, unless $v$ is unsafe, a \texttt{VDefVal} $ci'$ such that the step
is a \texttt{VEnv.AddDef} at every level.
\thesisref{axioms.tex, \S ``Definitions ($\delta$ reduction)''}
\end{lemma}
\begin{proof}
\uses{family:check-decls, family:add-const-wf, family:check-body, family:safe-primitives}
\leanok
Case split on $v.\mathtt{safety} = \mathtt{unsafe}$. The unsafe branch adds the header as an
axiom and then checks the body against that environment, concluding by \texttt{addUnsafeDef.WF};
the safe branch uses \texttt{checkDefinition.WF} and \texttt{addDef.WF}, with the
primitive-preservation side condition discharged from \texttt{hasPrimitives}.
\end{proof}

\begin{lemma}[Adding a theorem preserves well-formedness]
\label{thm:add-theorem-wf}
\lean{Lean4Lean.addTheorem.WF}
\leanok
\srcloc{Lean4Lean/Verify/Environment.lean}{75}
A successful $\mathtt{addTheorem}\ env\ v$ yields a well-formed model family extending the old
one, the step being a \texttt{VEnv.AddConst} of $\mathtt{thmInfo}\ v$ at every safety level: the
proof term is discarded, only its existence and the Prop-ness of the statement are recorded.
\thesisref{axioms.tex, \S ``Definitions ($\delta$ reduction)''}
\end{lemma}
\begin{proof}
\uses{family:check-decls, family:add-const-wf, ind:tr-env}
\leanok
\texttt{checkTheorem.WF} then \texttt{addConst.WF} with the \texttt{thm} step.
\end{proof}

\begin{lemma}[Adding an opaque definition preserves well-formedness]
\label{thm:add-opaque-wf}
\lean{Lean4Lean.addOpaque.WF}
\leanok
\srcloc{Lean4Lean/Verify/Environment.lean}{91}
A successful $\mathtt{addOpaque}\ env\ v$ gives a well-formed model family with a
\texttt{VEnv.AddConst} step for $\mathtt{opaqueInfo}\ v$ at each safety level; no defining
equation is registered, so the body is checked but unusable for $\delta$.
\thesisref{axioms.tex, \S ``Definitions ($\delta$ reduction)''}
\end{lemma}
\begin{proof}
\uses{family:check-decls, family:add-const-wf, ind:tr-env}
\leanok
\texttt{checkOpaque.WF} then \texttt{addConst.WF} with the \texttt{opaque} step.
\end{proof}

\begin{lemma}[Trivial \texttt{Except.WF} lemmas for the failure branches]
\label{family:except-wf-throw}
\lean{Lean4Lean.Except.WF.throw', Lean4Lean.Except.WF.throwBind}
\leanok
\srcloc{Lean4Lean/Verify/Environment.lean}{114}
Two private one-liners saying that a \texttt{throw}, possibly bound to a continuation, satisfies
any postcondition vacuously.
\end{lemma}
\begin{proof}
\leanok
There is no successful result to constrain.
\end{proof}

\begin{lemma}[Adding a mutual definition block preserves well-formedness]
\label{thm:add-mutual-wf}
\lean{Lean4Lean.addMutual.WF}
\leanok
\srcloc{Lean4Lean/Verify/Environment.lean}{120}
A successful $\mathtt{addMutual}\ env\ vs$ yields a well-formed model family extending the old one
at every safety level. The headers are checked first, as axioms, and the bodies against the
temporary environment holding all of them, which is what the \texttt{mutualDef} step expects.
\thesisref{axioms.tex, \S ``Definitions ($\delta$ reduction)''}
\end{lemma}
\begin{proof}
\uses{thm:add-mutual-block-wf, def:tr-mutual-header, family:check-constant-val, family:except-wf-throw}
\leanok
Two monadic loops verified through \texttt{TypeChecker.M.WF.forInFresh} and
\texttt{forInForall}$_2$, the first producing \texttt{TrMutualHeader} data for each member and the
second the translated bodies; concluded by \ref{thm:add-mutual-block-wf}.
\end{proof}

\begin{theorem}[Checked declaration addition preserves well-formedness]
\label{thm:add-decl-wf}
\contrib{mixed}
\lean{Lean4Lean.addDecl.WF}
\leanok
\uses{ind:tr-env, def:vtc-venvs}
\srcloc{Lean4Lean/Verify/Environment.lean}{197}
For any \texttt{Declaration}, a successful
$\mathtt{addDecl}\ env\ decl$ with checking enabled gives a well-formed model family $ves'$ that
extends $ves$ at every safety level. This is the top of the environment-side verification: it is
what would make \texttt{TrEnv} an invariant of kernel operation.
\thesisref{axioms.tex, \S ``Definitions ($\delta$ reduction)''}
\axfoot{sorryAx: the inductDecl case is itself a sorry}
\end{theorem}
\begin{proof}
\uses{thm:add-axiom-wf, thm:add-definition-wf, thm:add-theorem-wf, thm:add-opaque-wf, thm:add-mutual-wf, thm:add-quot-wf}
\textbf{Not proved in Lean.} One case per declaration kind, each discharged by the corresponding
lemma above; the \texttt{quotDecl} case is now discharged by the real \ref{thm:add-quot-wf}
contributed on \texttt{iota}, where master had a vacuous argument. The \texttt{inductDecl} case is
\texttt{sorry} (\srcloc{Lean4Lean/Verify/Environment.lean}{208}): what is missing is a
construction of an \texttt{AddInduct} witness from \texttt{Environment.addInductive}. The
\texttt{iota} contribution to this declaration itself is confined to the docstring, the signature
line and the \texttt{quotDecl} line.
\end{proof}

\section{The trust boundary}

\begin{definition}[\texttt{Std.TreeMap.all} and \texttt{any} agree with the list view]
\label{family:treemap-axioms}
\lean{Std.TreeMap.all_eq_all_toList, Std.TreeMap.any_eq_any_toList}
\leanok
\srcloc{Lean4Lean/Verify/Axioms.lean}{10}
Two axioms stating that \texttt{TreeMap.all} and \texttt{any} equal the corresponding
\texttt{List} predicates over \texttt{toList}. Both cite leanprover/lean4 issue 12798, that is,
they are missing upstream lemmas rather than unprovable facts.
\end{definition}

\begin{definition}[\texttt{PersistentArray} as a list]
\label{family:persistent-array-axioms}
\lean{Lean.PersistentArrayNode.toList', Lean.PersistentArray.WF, Lean.PersistentArray.toList', Lean.PersistentArray.toList'_empty, Lean.PersistentArray.toList'_push, Lean.PersistentArray.size_empty, Lean.PersistentArray.size_push, Lean.PersistentArray.WF.toList'_length}
\leanok
\srcloc{Lean4Lean/Verify/Axioms.lean}{22}
A noncomputable list view of \texttt{PersistentArray} with a \texttt{WF} predicate generated by
\texttt{empty} and \texttt{push}. Only \texttt{toList'\_push} is an axiom, because
\texttt{insertNewLeaf} is \texttt{partial}; the other facts are proved.
\end{definition}

\begin{definition}[\texttt{PersistentHashMap} as an association list]
\label{family:persistent-hash-map-axioms}
\lean{Lean.PersistentHashMap.Node.toList', Lean.PersistentHashMap.toList', Lean.PersistentHashMap.WF, Lean.PersistentHashMap.WF.toList'_insert, Lean.PersistentHashMap.WF.find?_eq, Lean.PersistentHashMap.findAux_isSome}
\leanok
\srcloc{Lean4Lean/Verify/Axioms.lean}{55}
The list view of \texttt{PersistentHashMap}, with three axioms: insertion permutes the list view,
\texttt{find?} is \texttt{lookup} on it, and \texttt{containsAux} agrees with \texttt{findAux}.
They are axioms because \texttt{insertAux}, \texttt{findAux} and \texttt{containsAux} are opaque.
Together with \texttt{sorryAx} these three are what
\texttt{\#print axioms Lean4Lean.TrEnv.proj\_defeq}
(\srcloc{Lean4Lean/Tests/ProjInhabit.lean}{595}) reports beyond \texttt{propext},
\texttt{Classical.choice} and \texttt{Quot.sound}, so they are on the critical path of the
contributed projection work.
\end{definition}

\begin{definition}[Structural equality of \texttt{Syntax}]
\label{family:syntax-axioms}
\lean{Lean.Syntax.structEq', Lean.Syntax.structEq'_node, Lean.Syntax.structEq_eq}
\leanok
\srcloc{Lean4Lean/Verify/Axioms.lean}{90}
A total copy \texttt{structEq'} of the \texttt{partial def Lean.Syntax.structEq}, an unfolding
lemma for the node case, and the axiom identifying the two.
\end{definition}

\begin{definition}[A total copy of \texttt{Lean.Level.normalize}]
\label{family:level-total}
\lean{Lean.Level.Total.size, Lean.Level.Total.normalize, Lean.Level.Total.mkMaxAux, Lean.Level.Total.skipExplicit, Lean.Level.Total.isExplicitSubsumedAux, Lean.Level.Total.getMaxArgsAux, Lean.Level.Total.one_le_size, Lean.Level.Total.size_getLevelOffset}
\leanok
\srcloc{Lean4Lean/Verify/Axioms.lean}{133}
\texttt{Lean.Level.normalize} and four helpers are \texttt{partial def}s and therefore opaque. The
\texttt{Total} namespace is a clause-by-clause copy under the same names with termination proofs
supplied, using \texttt{size} as the primary and a private \texttt{tag} as the secondary measure.
The section docstring pins the exact upstream file and line range (v4.33.0-rc2,
\texttt{Lean/Level.lean} lines 319 to 404) and points at a finite-corpus test; this is the model
of how such an assumption should be presented.
\end{definition}

\begin{definition}[Level normalization and level data]
\label{family:level-axioms}
\lean{Lean.Level.normalize_eq, Lean.Level.mkMaxAux_eq, Lean.Level.skipExplicit_eq, Lean.Level.isExplicitSubsumedAux_eq, Lean.Level.hasParam_eq, Lean.Level.hasMVar_eq, Lean.Level.instLawfulBEqLevel, Lean.Level.mkLevelIMaxCore_eq}
\leanok
\uses{family:level-total}
\srcloc{Lean4Lean/Verify/Axioms.lean}{177}
Eight axioms identifying the opaque \texttt{partial} level functions with their total copies, the
bit-packed \texttt{hasParam} and \texttt{hasMVar} accessors with structural definitions, the
C++-implemented \texttt{BEq Level} with a lawful instance, and a workaround for a known
\texttt{mkLevelIMaxCore} elaboration issue.
\thesisref{unique.tex and axioms.tex: the level inequality algorithm}
\end{definition}

\begin{definition}[The C++-implemented \texttt{Expr} primitives]
\label{family:expr-axioms}
\lean{Lean.Expr.looseBVarRange_eq, Lean.Expr.replace_eq, Lean.Expr.liftLooseBVars_eq, Lean.Expr.lowerLooseBVars_eq, Lean.Expr.instantiate1_eq, Lean.Expr.instantiate_eq, Lean.Expr.instantiateRev_eq, Lean.Expr.instantiateRange_eq, Lean.Expr.instantiateRevRange_eq, Lean.Expr.abstract_eq, Lean.Expr.abstractRange_eq, Lean.Expr.hasLooseBVar_eq, Lean.Expr.eqv_eq, Lean.Expr.equal_eq}
\leanok
\srcloc{Lean4Lean/Verify/Axioms.lean}{360}
Fourteen axioms identifying Lean's C++ implementations of de Bruijn lifting, lowering,
instantiation, abstraction, loose-bvar queries and structural equality with reference Lean
definitions (\texttt{liftLooseBVars'}, \texttt{instantiate1'}, \texttt{abstractList},
\texttt{eqv'} and so on). Most are annotated as candidates for \texttt{@[implemented\_by]}, that
is, performance shims rather than genuine gaps.
\reviewnote{This is the largest single trust assumption of the project: every kernel-level theorem
about \texttt{Expr} manipulation rests on it. The identifications are shallow and syntactic, but
nothing in the repository checks them beyond inspection.}
\end{definition}

\section{Contribution summary}

The \texttt{iota} branch contributed the inductive-block half of this layer. On master
\texttt{AddInduct} was a \texttt{Prop} with no constructors, documented as ``essentially a
\texttt{sorry}'': the \texttt{induct} step of \ref{ind:tr-env} could never fire, so every
environment containing an inductive type lay outside the verified relation, and every consequence
of $\mathtt{TrEnv}'$ was silently restricted to inductive-free environments. The branch replaced it
with a real witness (\ref{struct:add-induct}) recording the kernel data, the four stages of
\texttt{VEnv.addInduct}, a \ref{struct:tr-ind-type} per type former and a \ref{struct:tr-recursor}
per recursor, freshness of every inserted name, and the kernel's actual insertion order as a
permutation of the model's stage order (\ref{def:add-induct-consts}). Supporting it is the
\texttt{insertConsts} map calculus (\ref{def:insert-consts}, \ref{family:insert-consts-lemmas}).
From the witness the branch \emph{derives} the block's bookkeeping instead of assuming it as
further fields: stage recomposition (\ref{family:add-induct-stages}), membership by kind and
absence of $\delta$-values (\ref{family:add-induct-membership}), agreement of the kernel names with
the model's together with their distinctness (\ref{family:add-induct-names}), map well-formedness
(\ref{thm:add-induct-wf}) and two-way lookup transport (\ref{family:add-induct-find}).

The three interface theorems \ref{thm:add-induct-rec-find}, \ref{thm:add-induct-rec-reg} and
\ref{thm:add-induct-ctor-find} are what the rest of the chapter consumes. They feed, on the
alignment side, \ref{thm:aligned-add-rules} and \ref{thm:aligned-add-induct}, which replace
master's one-line \texttt{nomatch} for the \texttt{induct} case and isolate the whole
kernel-order-versus-model-order difficulty in a single permutation argument; and, on the reduction
side, \ref{thm:tr-env-pats-iota}, which says every kernel recursor rule is a registered $\iota$
pattern of the model, and its composition \ref{thm:tr-env-iota-rec} with
\ref{thm:tr-env-iota-defeq}. A second, unplanned consequence is the whole of
\srcloc{Lean4Lean/Verify/Environment/Quot.lean}{1}: on master \texttt{checkEqType.WF} proved
\texttt{False}, because a \texttt{TrEnv} could not contain the inductive \texttt{Eq}, so quotient
initialization was well-formed by vacuity. Once \texttt{AddInduct} became inhabitable that argument
evaporated and the branch had to prove it for real, which is \ref{thm:with-local-decl-run} through
\ref{thm:add-quot-wf}: a symbolic replay of the kernel's expression builder
(\ref{def:quot-telescope-data}, \ref{family:quot-telescope-find}, \ref{def:quot-tactics}), the
evaluation of the four \texttt{mkForall} telescopes (\ref{family:quot-type-eqs},
\ref{thm:add-quot-eq}), the analysis of \texttt{checkEqType} (\ref{def:check-eq-type-shape},
\ref{thm:environment-get-ok}, \ref{thm:check-eq-type-ok}), explicit translations of the four
quotient types (\ref{family:quot-type-tr}, \ref{thm:quot-ready}), and the assembly
\ref{thm:exists-add-quot}, \ref{thm:add-quot-wf}. That last theorem discharges the
\texttt{quotDecl} case of \ref{thm:add-decl-wf} for real.

The \texttt{trproj} branch contributed the structure-likeness and projection half. Its additions to
the inductive-block vocabulary are the parameters and the \texttt{all}, \texttt{numParams},
\texttt{numIndices} and arity fields of \ref{struct:tr-ind-type} and the \texttt{name\_major} field
of \ref{struct:tr-recursor}. On them rest \ref{thm:aligned-constants-pull},
\ref{thm:tr-env-ctor-arity} (a constructor's model type has $\Pi$-arity
$numParams + numFields$) and \ref{thm:tr-env-structure-rec}, which takes exactly the three tests
the kernel's \texttt{inferProj} and \texttt{reduceProjCore} perform before accepting a projection
and returns the recursor's telescope split and the constructor's parameter count. The branch then
built the inverse of the $\iota$ interface: \ref{struct:iota-rule} bundles the kernel data behind
one registered pattern together with its model-side shape, \ref{thm:iota-rule-step} transports it
along the non-\texttt{induct} steps, and \ref{thm:pats-iota-inv-shape} (restated at the kernel
environment level as \ref{thm:tr-env-pats-iota-inv-shape}) recovers a kernel rule from a registered
pattern. With the $\beta$-telescope calculus \ref{family:beta-telescope} these combine into
\ref{thm:tr-env-proj-defeq}: the recursor expansion of the $i$-th projection of a value
definitionally equal to a saturated constructor spine is definitionally equal to the $i$-th field.
The architecture is the right one, in that the projection computation rule is derived from the
registered $\iota$ rule rather than postulated alongside it.

Connection to the rest of the blueprint. The $\iota$ interface reaches the kernel model at exactly
one point, \ref{thm:vtc-iota-reduce-core}, which has no \texttt{sorry} of its own; the projection work reaches it at one
point through \ref{thm:tr-expr-mk-app-list}, and otherwise not at all. Upstream, the two branches
depend on the inductive-specification lemmas of Chapter~\ref{chap:inductive} and on the projection
translation \ref{struct:trprojctor} of Chapter~\ref{chap:trexpr}; downstream, the intended
consumers are \ref{thm:vtc-reduceprojcore} and \ref{thm:vtc-inferproj-struct} in
Chapter~\ref{chap:typechecker}.

What remains open. Nothing in the repository constructs an \texttt{AddInduct} witness from
\texttt{Environment.addInductive}, and the \texttt{inductDecl} case of \ref{thm:add-decl-wf} is
still \texttt{sorry}, so every theorem above is conditional on a hypothesis that is never
discharged: the branches moved the \texttt{sorry} from ``\texttt{AddInduct} is empty'' to
``\texttt{AddInduct} has no producer''. \ref{thm:tr-env-proj-defeq} has no consumer outside
\ref{test:projinhabit-axioms}, because \ref{thm:vtc-reduceprojcore}, \ref{thm:vtc-inferproj-struct}
and \ref{thm:vtc-inferproj} are \texttt{sorry} (\ref{thm:vtc-reduceproj}, which delegates to the
first of these, is itself proved) and its premise \texttt{TrProjCtor} is so far inhabited
only by hand in the tests. Finally both \ref{family:beta-telescope} and
\ref{thm:tr-env-proj-defeq} inherit \texttt{sorryAx} through unique typing
(\texttt{VEnv.WF.orderedStrong}, \texttt{VEnv.WF.patsStrong}), which is outside this group but on
the critical path of its headline result.

\section{Review notes}

\begin{itemize}
\item \textbf{High.} The \texttt{inductDecl} case of \texttt{addDecl.WF} is still \texttt{sorry}
(\srcloc{Lean4Lean/Verify/Environment.lean}{208}) and nothing anywhere constructs an
\texttt{AddInduct} witness from \texttt{Environment.addInductive}. The entire inductive-block
apparatus of both branches (\ref{struct:add-induct}, \ref{struct:tr-ind-type},
\ref{struct:tr-recursor}, \ref{thm:aligned-add-induct}, \ref{thm:tr-env-pats-iota},
\ref{thm:pats-iota-inv-shape}) is conditional on that hypothesis. The theorems are correct but in
the end-to-end sense currently vacuous.
\item \textbf{High.} \texttt{TrEnv.proj\_defeq} is orphaned
(\srcloc{Lean4Lean/Verify/Environment/Lemmas.lean}{1021}): its only consumer in the repository is
the \texttt{\#print axioms} check at \srcloc{Lean4Lean/Tests/ProjInhabit.lean}{595}.
\texttt{reduceProjCore.WF} (\srcloc{Lean4Lean/Verify/TypeChecker/Reduce.lean}{143}),
\texttt{inferProj.WF\_struct} (\srcloc{Lean4Lean/Verify/TypeChecker/InferType.lean}{392}) and
\texttt{inferProj.WF} (\srcloc{Lean4Lean/Verify/TypeChecker/InferType.lean}{407}) remain
\texttt{sorry}, so the strongest \texttt{trproj} theorem does not yet connect to the kernel
model. Its premise \texttt{TrProjCtor}
is likewise only ever inhabited by hand in tests, the general construction
\texttt{inferProj.WF\_struct} being itself \texttt{sorry}.
\item \textbf{Medium.} \texttt{TrEnv.pats\_iota\_inv\_shape}
(\srcloc{Lean4Lean/Verify/Environment/Lemmas.lean}{903}) is redundant: its proof is literally
\texttt{TrEnv'.pats\_iota\_inv\_shape H hp}, the two statements being definitionally equal, and
unlike its sibling \texttt{TrEnv.pats\_iota'} it does not convert \texttt{SMap.find?} into
\texttt{Environment.find?}. Its docstring claims otherwise, and \texttt{proj\_defeq} still performs
the conversions by hand at \srcloc{Lean4Lean/Verify/Environment/Lemmas.lean}{1060} and
\srcloc{Lean4Lean/Verify/Environment/Lemmas.lean}{1072}.
\item \textbf{Medium.} \texttt{TrRecursor.all}
(\srcloc{Lean4Lean/Verify/Environment/Basic.lean}{240}) and \texttt{TrRecursor.k}
(\srcloc{Lean4Lean/Verify/Environment/Basic.lean}{245}) are never used anywhere in the repository.
They are extra proof obligations on a producer of \texttt{AddInduct} that buy nothing; \texttt{k}
in particular records the K-like reduction flag, for which the model has no corresponding rule.
\item \textbf{Medium.} \texttt{insertConsts\_find?\_none}
(\srcloc{Lean4Lean/Verify/Environment/Basic.lean}{165}) is dead code: apart from its own statement,
its only occurrence is its own recursive call.
\item \textbf{Medium.} \texttt{AddInduct} is declared without a \texttt{: Prop} ascription
(\srcloc{Lean4Lean/Verify/Environment/Basic.lean}{272}), so it is a \texttt{Type}-valued structure
and the \texttt{induct} constructor of $\mathtt{TrEnv}'$ carries data. This is necessary, since the
derived lemmas project the fields \texttt{ivals}, \texttt{envR} and \texttt{order}, and harmless
for the present consumers, which all eliminate into \texttt{Prop}; but it is a real change to the
shape of the project's central invariant relative to master's
\texttt{inductive AddInduct ... : Prop} and it is nowhere documented. (It does not, as one might
read it, cost $\mathtt{TrEnv}'$ large elimination: with nine constructors it never had any.)
\item \textbf{Medium.} Misfiled material. The $\beta$-telescope family
(\srcloc{Lean4Lean/Verify/Environment/Lemmas.lean}{959}) is pure \texttt{VExpr} metatheory with no
environment-translation content and belongs in \texttt{Theory/Typing/Lemmas.lean}; the same holds
for \texttt{TrExpr.mkAppList} (\srcloc{Lean4Lean/Verify/Environment/Lemmas.lean}{1166}), which is
\texttt{TrExpr} rebuilding rather than environment bookkeeping. Both live here only because
\texttt{proj\_defeq} needs them.
\item \textbf{Medium.} \texttt{TrEnv.iota\_defeq}
(\srcloc{Lean4Lean/Verify/Environment/Lemmas.lean}{917}) is a three-line wrapper over
\texttt{VEnv.IsDefEq.pat} with no \texttt{TrEnv} hypothesis at all, yet is named as if it had one.
\item \textbf{Medium.} Brittleness in \texttt{Quot.lean}: \texttt{AddQuotAux}
(\srcloc{Lean4Lean/Verify/Environment/Quot.lean}{27}) hand-replays the binder names, binder infos
and \texttt{withLocalDecl} order of \texttt{Lean4Lean/Quot.lean}, with nothing linking the two
beyond the proofs breaking, and the \texttt{T\ldots\_tr} derivations
(\srcloc{Lean4Lean/Verify/Environment/Quot.lean}{331}) are sixty lines of hand-built
\texttt{TrExprS} trees. Correct, but expensive to maintain.
\item \textbf{Low.} \texttt{addQuot.WF} extracts the four \texttt{checkName} results with four
near-identical seven-line blocks (\srcloc{Lean4Lean/Verify/Environment/Quot.lean}{521}) and
rebuilds the \texttt{safePrimitives} field with a 22-line hand-nested chain of
\texttt{Environment.find?\_add\_of\_ne} (\srcloc{Lean4Lean/Verify/Environment/Quot.lean}{584});
both are mechanical repetition that a single auxiliary lemma would remove. Also
\texttt{Lean4Lean.withLocalDecl\_run} (\srcloc{Lean4Lean/Verify/Environment/Quot.lean}{20}) and
\texttt{Lean4Lean.AddQuotAux.Environment.get\_ok}
(\srcloc{Lean4Lean/Verify/Environment/Quot.lean}{244}) are generally useful lemmas placed in,
respectively, the root namespace and a private auxiliary namespace.
\item \textbf{Low (reviewer caveat, not a defect).} Some lines the per-line attribution marks as
contributed are master code relocated by the refactor commits:
\texttt{TrConstant.sf\_mono}/\texttt{mono}, \texttt{TrConstVal.mono} and \texttt{TrDefVal.mono}
(\srcloc{Lean4Lean/Verify/Environment/Basic.lean}{34}, moved from \texttt{Lemmas.lean}) and
\texttt{insertDefs\_wf} (\srcloc{Lean4Lean/Verify/Environment/Lemmas.lean}{211}, moved from
\texttt{Extension.lean}). Sixteen of the 398 \texttt{iota} lines in \texttt{Basic.lean} are of
this kind, so the raw line counts overstate the branch slightly.
\end{itemize}
